\section {Dantzig Wolfe Decomposition (DWD) of G-LRSP} 
\label{part1:dantzigW}
In this section, we prove that the SPP-LRSP 
can be obtained by applying Dantzig Wolfe decomposition (DWD) to the graph-based 
formulation of the LRSP (G-LRSP). 

\subsection{DWD in General}

The DWD method is used in integer programming to improve  the LP relaxation bound. 
The method reformulates the original problem in such a way that the reformulation 
yields a tighter bound for the optimal objective function value. To explain the 
steps of the DWD method in generality, suppose we have an integer program, P that 
can be partitioned into two sets of constraints. 
\begin{optprog}
   Minimize & \objective{c^Tx} \nonumber\\
(P) \ \  subject to & A^1x & \geq & d^1, \label{sys1}\\
           & A^2x & \geq & d^2, \label{sys2} \\
           & x & \geq & 0,\mbox{ and integer}  \label{sys3}     
\end{optprog} 
where $A^1 \in \mathbb{Q}^{m\times n}$, and $A^2 \in \mathbb{Q}^{l \times n}$ are
constraint matrices and $c \in \mathbb{Q}^{n}$, $d^1 \in \mathbb{Q}^{m}$
and $d^2 \in \mathbb{Q}^{l}$ are rational vectors. 
By relaxing the integrality constraint on $x$, an LP relaxation bound is obtained.
The idea of DWD method is to get possibly a better LP relaxation bound by 
optimizing the intersection of constraints (\ref{sys1}) with the 
convex hull of (\ref{sys2}) and (\ref{sys3}). 
%$U \in \mathbb{R}_+^n$ is the variable vector and the integrality restrictions 
%for $U$ are included in constraints (\ref{sys2}). 
Let $X^2=\{x \in \mathbb{Z}^n:A^2x \geq d^2, x \geq 0\}$ be the set of integer points 
defined by systems (\ref{sys2}) and (\ref{sys3}), and $Conv(X^2)$ be the convex hull of $X^2$. 
Assuming that the $Conv(X^2)$ is bounded, we can express $Conv(X^2)$ using the convex 
combination of the extreme points of $X^2$:
\begin{optprog} 
& Conv(X^2)=\{x \in \mathbb{R}^n | x = \sum_{q \in E}\mbox{x}^q\theta_q, \sum_{q \in E}\theta_q=1, 
\theta \in \mathbb{R}^{|E|}_+\} \label{convhull}
\end{optprog} 
where $\{\mbox{x}^q\}_{q \in E}$ is the set of extreme points of $X^2$. Here, we assume 
that $Conv(X^2)$ is bounded and pursue the review of DWD method under this assumption since 
this is the case in LRSP.  For unbounded case, the definition of $Conv(X^2)$ in (\ref{convhull}) 
can be updated by considering the extreme rays of $X^2$.  

\zaupdate{We can use decomposable subproblem as an example, like LRSP?}

%To find a bound for the optimal objective function value of P, the method 
%optimizes the objective function over the intersection of system (\ref{sys1}) and 
%$Conv(X)$. 
Then, LP relaxation of P can be reformulated using (\ref{convhull}):
\begin{optprog}
 Minimize & \objective{c(\sum_{q \in E}\mbox{x}^q\theta_q)} \nonumber\\
(P$_{\mbox{M}}$) \ \ subject to & A^1(\sum_{q \in E}\mbox{x}^q\theta_q ) & \geq & d^1,  \nonumber \\
           & \sum_{q \in E}\theta_q & = & 1,  \nonumber\\
           & \theta & \in & \mathbb{R}^{|E|}_+.\nonumber
          % & \theta & \in & \mathbb{I}^{|R|}_+ \label{integerlambda}
\end{optprog} 
If we rearrange,
\begin{optprog}
 Minimize & \objective{\sum_{q \in E}(c\mbox{x}^q)\theta_q} \nonumber\\
(P$_{\mbox{M}}$) \ \ subject to & \sum_{q \in E}(A^1\mbox{x}^q)\theta_q  & \geq & d^1, & (\pi)\label{master-con1}\\
           & \sum_{q \in E}\theta_q & = & 1, & (\mu) \label{master-con2}\\
           & \theta & \in & \mathbb{R}^{|E|}_+.  \label{binarylambda}
          % & \theta & \in & \mathbb{I}^{|R|}_+ \label{integerlambda}
\end{optprog} 
%Since P is an integer linear program we have to force the integrality of variables.
%To do that, we only consider the feasible integer points of $Conv(X)$
%using constraints (\ref{binarylambda}) and (\ref{integerlambda}). 
P$_{\mbox{M}}$ is called {\em master} problem. 
The bound obtained by solving P$_{\mbox{M}}$ is typically tighter than the bound
obtained by solving the LP relaxation of P. If system (\ref{sys2}) has integrality  
property, i.e., all extreme points of system (\ref{sys2}) are integer, then the LP
relaxation value of P is equal to the objective function value of P$_{\mbox{M}}$. 
Otherwise, solution of P$_{\mbox{M}}$ yields a better lower bound to the
optimal solution value of P. 

%There is one variable corresponds to an extreme point in P$_{\mbox{M}}$. 
To find the solution of  P$_{\mbox{M}}$, we have to explore the set of extreme points, 
$\{\mbox{x}^q\}_{q \in E}$. One way of doing that is to enumerate the whole set.
However, in most cases enumeration of the whole set is not possible because of the 
size of set $E$. Other way of finding the solution of P$_{\mbox{M}}$ is 
to use {\em column generation algorithm}. Column generation is an approach for 
solving LPs for which we want to avoid explicitly enumerating all columns. 
By solving a sequence of linear programs and dynamically generating columns eligible to
enter the basis, such an algorithm implicitly considers all columns but generates
only those that may improve the objective function.
The basic idea is to determine a subset of the variables that includes those 
that are nonzero in some optimal solution. By solving the LP including only 
this subset of columns rather than all possible columns, we can obtain the 
optimal LP solution.

In each iteration of the column generation algorithm, a restricted version of 
P$_{\mbox{M}}$ (RMP) that includes only a subset of columns is solved.
The column generation algorithm feeds the dual information from the solution of the 
RMP to an optimization model called {\em pricing} problem (or {\em subproblem}). 
The pricing problem searches columns in set $X^2$ with an objective function 
equal to the reduced cost of a variable in P$_{\mbox{M}}$. 
The pricing problem can be formulated as follows:
\begin{optprog}
Minimize & \objective{z_{rc} = (c-\pi A^1)x - \mu} \nonumber \\
subject to &  A^2x & \geq & d^2 \nonumber\\
            & x & \geq & 0,\mbox{ and integer}  \nonumber
\end{optprog}
where $\pi$ is the dual vector corresponds to system (\ref{master-con1})
and $\mu$ is the dual variable for constraint (\ref{master-con2}). 
%and $rc$ is the reduced cost of a variable in RMP.
The objective of the pricing problem is to either identify new columns to add to 
the RMP or prove optimality of the current solution. From the theory of linear 
programming, we know that if every variable has a non-negative reduced cost, 
then the solution is optimal. If $z_{rc}<0$ for a given $(\pi, \mu)$, then a column 
in $\{\mbox{x}^q\}_{q \in E}$ is found and can be added to the RMP. Then, the RMP is 
reoptimized. The algorithm continues until $z_{rc}\geq 0$ for a given dual vector. 

%If the pricing problem identifies at least one pairing with negative reduced cost, then we add
%one or more columns to the RMP and reoptimize. Otherwise, the current solution is optimal for
%the full problem.

With column generation algorithm, we can find a lower bound for the optimal integer 
solution value for P. To find the optimal integer solution using model P$_{\mbox{M}}$, 
we have to force the integrality of variables. This can be done by adding integrality 
restriction to variable $\theta$ in  P$_{\mbox{M}}$. 

\zaupdate{Mention branch and price and say will be explained how to solve 
integer master problem?}


\subsection{DWD of G-LRSP}

Back to the G-LRSP, to find the DWD of the problem, 
we decompose the problem into two subsystems: master problem and subproblem.
The {\em master problem} is defined by constraints (\ref{set-part}) and (\ref{FacCap}),
while the {\em subproblem} is defined by constraints (\ref{enterleave}), (\ref{VehCap}), 
(\ref{flow-conv}), (\ref{timeLim}), (\ref{fix-0}), (\ref{x-var}), (\ref{y-var}), 
(\ref{t-var}), and (\ref{dri-var}).
Remember that the vehicle index set, $H$ is 
equal to the union of the vehicle indices located at each facility, i.e., 
$H = \cup_{j \in J}H_j$. Thus, all constraints of the subproblem are 
indexed by the set J and it is therefore immediate
that the subproblem decomposes by facility. 

For each candidate facility $j \in J$, the set of integer solutions of the 
decomposed subproblem can be represented by the set of vectors defined as
\begin{eqnarray}
& & \{ \ \  (x, t, v)^j \in \Z^{(|N|\times |N| \times |H_j|) \times \{0, 1\} \times(|H_j|)}, 
 (y)^j \in \R^{(|N|\times |N| \times |H_j|)}  \mid  \ \  (x, y, t, v)^j \mbox{ satisfies } \nonumber\\
& &  \ \ (\ref{enterleave})_j, 
 (\ref{VehCap})_j, (\ref{flow-conv})_j, (\ref{timeLim})_j, (\ref{fix-0})_j, 
(\ref{x-var})_j, (\ref{y-var})_j, (\ref{t-var})_j, (\ref{dri-var})_j \}, \label{subsys} 
\end{eqnarray}

which is described only by those constraints associated with facility
$j$ and specifies values only for those variables associated with
vehicles assigned to facility $j$. Hence, the index $j$ for the
constraints represents the set of constraints associated with facility
$j$. System (\ref{subsys}) is bounded for all facilities.
We let $E$ be a set indexing the extreme points %members 
of all of the above sets,
with $E_j$ the indices for vectors associated with facility $j$ only,
so that $E = \cup_{j \in J} E_j$. For a facility $j$ and an index $q
\in E_j$, the corresponding member $(\mbox{x}^q, \mbox{y}^q,
\mbox{t}^q,\mbox{v}^q)$ of the above set is then a vector with the following
interpretation: 
\begin{eqnarray*}
\mbox{x}_{ikh}^q & = & \left \{ \begin{array}{ll}
                 1  & \mbox{if vehicle $h$ travels on arc $(i,k)$ in solution 
		   $q$, $\forall (i, k) \in A$, and $h \in H_j$}\\
                 0  & \mbox{otherwise}                                                  
                \end{array}
       \right. \\
\mbox{y}_{ikh}^q & = & \mbox{Load of vehicle $h$ on arc $(i,k)$ in solution 
		   $q$, $\forall (i, k) \in A$, and $h \in H_j$ }\\ 
\mbox{t}^q & = & \left \{ \begin{array}{ll}
                 1  & \mbox{if facility $j$ is open in solution $q$ }\\
                 0  & \mbox{otherwise}                                                  
                \end{array}
       \right.\\
\mbox{v}_{h}^q & = & \left \{ \begin{array}{ll}
                 1  & \mbox{if vehicle $h$ is used in solution $q$, $\forall h \in H_j$}\\
                 0  & \mbox{otherwise}                                                  
                \end{array}
       \right. 
\end{eqnarray*}
%ALTERNATE TO THE PREVIOUS FORMAT%%%%%%%%%%%%%%%%%%%%%%%%%%%%%%%%%%%%%%%%%%%%%%%%%%%%%%%
%for each pair $(i, k) \in A$ and $h \in H_j$, the parameter 
%$\mbox{x}_{ikh}^q = 1$ if vehicle $h$ travels on arc $(i,k)$ in solution 
%$q$ and is 0 otherwise; $\mbox{y}_{ikh}^q$ is equal to the amount of 
%demand carried by vehicle $h$ on arc $(i,k)$; $\mbox{t}^q =1$ if facility 
%$j$ is open in solution $q$ and is 0 otherwise; and for each $h \in H_j$, 
%$\mbox{v}_{h}^q = 1$ if vehicle $h$ is used in solution $q$ and is 0 otherwise. 
%%%%%%%%%%%%%%%%%%%%%%%%%%%%%%%%%%%%%%%%%%%%%%%%%%%%%%%%%%%%%%%%%%%%%%%%%%%%%%%%%%%%%%%%%

Note that the variable $t$ indicating whether the facility is open does not appear in any 
of the linear constraints of the subproblem and can hence be set to either 0
or 1 without affecting feasibility. 

Following the DWD methodology, we can reformulate master problem using the convex 
combination of the extreme points of the subproblem. 
Because the subproblem decomposes as described above, solutions to the
original problem can be seen as vectors obtained by ``recomposing'' convex
combinations of the members of $E_j$ for each $j \in J$. In other words, any
solution $(x, y, t, v)$ to the master problem can be
written as:
\begin{optprog}
& x_{ikh} & = & \sum_{q \in E}\mbox{x}_{ikh}^q \theta_q &  \forall  (i,k) \in A, h \in H,  \label{XintermsTeta1}\\
& y_{ikh} & = &  \sum_{q \in E}\mbox{y}_{ikh}^q \theta_q & \forall  (i,k) \in A, h \in H, \label{TotalYintermsTeta1}\\
& t_j  & = & \sum_{q \in E_j}\mbox{t}^q\theta_q   &  \forall  j \in J, \label{Tvartrans1}\\
& v_h  & = & \sum_{q \in E}v_{h}^q\theta_q &  \forall h \in H, \label{VintermsTeta1}\\
& \sum_{q \in E_j}\theta_q & = & 1 & \forall  j \in J, \label{eq11}\\
& \theta_q  & \geq & 0 &  \forall q \in E. \label {conv1}
\end{optprog}  

\zaupdate{change it to $E_j$?}

Using (\ref{XintermsTeta1}) - (\ref {conv1}), we can then formulate the
master problem as:
\begin{optprog}
{\bf (MDW)} Minimize & \objective{\sum_{q \in E} \tilde{C_{q}} \theta_q} \nonumber\\
s.t. & \sum_{q \in E} b_{iq} \theta_q & = & 1 & \forall i \in I, \label{spar1} \\
           & -\sum_{q \in E_j} \sum_{i \in I} b_{iq} D_i \theta_q  +
L^F_j \sum_ {q \in E_j}\mbox{t}^q \theta_q & \geq & 0 & \forall j \in J, \label{rel-2a}  \\
           & \sum_{q \in E_j} \theta_q  & = & 1 & \forall j \in J,\label{conv2a} \\
           & \theta_q  & \geq & 0 &  \forall q \in E, \label {conv3a}
\end{optprog} 
where 
%\begin{eqnarray*}
\begin{optprog}
& b_{iq} & = & \sum_{h \in H}\sum_{k \in N}\mbox{x}_{kih}^q, & \forall i \in I, q \in E, \nonumber\\
& \tilde{C_{q}} & = & C^F_q {\mbox{t}}^q+ C^V\sum_{h \in H}\mbox{v}_{h}^q + C^O\sum_{i \in N} \sum_{k \in N}\sum_{h \in H}T_{ik}\mbox{x}_{ikh}^q, 
& \forall  q \in E, \nonumber 
%\end{eqnarray*}
\end{optprog}
where $C^F_q = C^F_j$ when $q \in E_j$ for $j \in J$. Here, $ b_{iq}$ can be
interpreted as the number of times customer $i$ is visited in solution
$q$, and $\tilde{C_{q}}$ is the cost of solution $q$ (including
facility fixed cost and vehicle fixed costs) for all $q\in E$. 

%%%%%%%%%%%%%%%%%%%%%%%%%%%%%% COMMENT %%%%%%%%%%%%%%%%%%%%%%%%%%%%%%%%%%%%%%%%%%%%%%%%%
\begin{comment}
System (\ref{sys2}) is bounded, therefore set of extreme rays is empty and we can reformulate (EB-LRS) using
the extreme points of system (\ref{sys2}). Let $E$ be the extreme points of system (\ref{sys2}).
%$S=\{X'|BX' & \leq & D_2\}$
Define following parameters to represent a solution $q = (x_q^T, y_q^T, v_q^T, t_q^T )^T \in E$.
\begin{eqnarray*}
x_{ikhq} & = & \left \{ \begin{array}{ll}
                 1  & \mbox{if vehicle $h$ visits arc $(i,k)$ (i.e., $X_{ikh}=1$) in $q\in E$, $\forall i, k \in I, h \in H$ }\\
                 0  & \mbox{otherwise}                                                  
                \end{array}
       \right. \\
y_{ikhq} & = & \mbox{Load of vehicle $h$ on arc $(i,k)$ in $q\in E$, $\forall i, k \in I, h \in H$ }\\ 
v_{hq} & = & \left \{ \begin{array}{ll}
                 1  & \mbox{if vehicle $h \in H$ is used (i.e., $V_h=1$) in solution $q\in E$, $\forall h \in H$}\\
                 0  & \mbox{otherwise}                                                  
                \end{array}
       \right. \\ 
t_{jq} & = & \left \{ \begin{array}{ll}
                 1  & \mbox{if facility $j$ is open in solution $q\in E$, i.e., $T_j=1$, $\forall j \in J$ }\\
                 0  & \mbox{otherwise}                                                  
                \end{array}
       \right.
\end{eqnarray*}
To be able to simplify the formulation of master problem, we define following parameters.
\begin{eqnarray*}
g_{ijq} & = & \sum_{h \in H_j}\sum_{k \in I}x_{ikhq}, \hspace{0.2in} \forall i \in I, j \in J, q \in E \label{transform-capcon} \\
\mbox{}  & = & \mbox{number of times customer $i$ is served from facility $j$ in solution $q \in E$} \\
 b_{iq} & = & \sum_{h \in H}\sum_{k \in I}x_{ikhq} =  \sum_{j \in J}g_{ijq}, \hspace{0.2in} \forall i \in I, q \in E \label{transform-xa}\\
\mbox{} & = & \mbox{number of times customer $i$ is visited in solution $q \in E$} \\
\tilde{C_{q}} & = & \sum_{j \in J}C^F_j t_{jq}+ VF\sum_{h \in H}v_{hq}+OC \sum_{i \in I} \sum_{k \in I}\sum_{h \in H} TD_{ik}x_{ikhq}, \hspace{0.2in} \forall  q \in E \label{transform-cost} \\
\mbox{} & = & \mbox{cost of solution $q \in E$ including facility and vehicle fixed costs and the operating cost of vehicles.} 
\end{eqnarray*}

Using these parameters, master problem (Master-DW-1) is formulated as follows:
\begin{optprog}
{\bf (Master-DW-1)} \nonumber\\
Minimize & \objective{\sum_{q \in E} \tilde{C_{q}} \theta_q} \nonumber\\
s.t. & \sum_{q \in E} b_{iq} \theta_q & = & 1 & \forall i \in I \label{spar} \\
           & \sum_{q \in E} \sum_{i \in I} g_{ijq} D_i \theta_q  - CapF_{j} \sum_ {q \in E}t_{jq}\theta_q & \leq & 0 & \forall j \in J \label{rel-2}  \\
           & \sum_{q \in E} \theta_q  & = & 1 \label{conv2} \\
           & \theta_q  & \in & \{0,1\} &  \forall q \in E \label {conv3}
\end{optprog} 
where
\begin{eqnarray*}
\theta_q & = & \left \{ \begin{array}{ll}
                 1  & \mbox{if solution $q$ is selected, $\forall q \in E$}\\
                 0  & \mbox{otherwise}                                                  
                \end{array}
       \right. 
\end{eqnarray*}
We can retrieve original variables using the following relations:
\begin{optprog}
& X_{ikh} & = & \sum_{q \in E}x_{ikhq} \theta_q &  \forall  i,k \in I,  j \in J, h \in H \label{XintermsTeta}\\
& Y_{ikh} & = &  \sum_{q \in E}y_{ikhq} \theta_q & \forall  i,k \in I,  j \in J, h \in H \label{TotalYintermsTeta}\\
& T_j & =& \sum_{q \in E}t_{jq}\theta_q &  \forall  j \in J \label{Tvartrans}\\
& V_h & = & \sum_{q \in E}v_{hq}\theta_q &  \forall h \in H 
\end{optprog}  
\end{comment}
%%%%%%%%%%%%%%%%%%%%%%%%%%%%%%%%%%%%%%%%%%%%%%%%%%%%%%%%%%%%%%%%%%%%%%%%%%%%%%%%%%%%%%%%%%%%%%%%

To solve MDW, we can use column generation algorithm. 
Let $(\pi ^ T, \mu ^T, \alpha^T)^T$ be a dual vector such that
$\pi$ is the dual variable associated with constraint (\ref{spar1}), 
$\mu$ is the dual variable associated with (\ref{rel-2a}), 
and $\alpha$ is the dual variable associated with (\ref{conv2a}). 
The reduced cost of a column $q \in E_j$ for some $j \in J$ 
corresponds to variable $\theta_q$ is 
\begin{equation}
\label{dw-red-cost}
\hat{C}_{q} = \tilde{C_{q}} - \sum_{i\in I}b_{iq}\pi_{i} + 
\left(\sum_{i \in I}b_{iq}D_{i} - L^F_j \mbox{t}^{q}\right)\mu_{j} - \alpha_j
\end{equation}
Using the definitions of $\tilde{C}$,  we can rewrite (\ref{dw-red-cost}).
\begin{eqnarray}
\label{dw-red-cost-open}
\hat{C}_q & = &(C^F_j-L^F_j \mu_j)\mbox{t}^q + C^V \sum_{h \in H_j}\mbox{v}_h^q - \nonumber\\
              & & \sum_{h \in H_j}\sum_{k \in N} (\sum_{i \in I} (C^O T_{ki} - \pi_{i} 
              + \mu_{j}D_{i})\mbox{x}_{kih}^q + C^OT_{kj}\mbox{x}_{kjh}^q)- \alpha_j
\end{eqnarray} 
The column generation algorithm searches for the columns with most negative 
reduced cost defined by (\ref{dw-red-cost-open}) over the set of extreme points
for facility $j$ defined in (\ref{subsys}). This problem is called subproblem  
for facility $j \in J$, and can be formulated explicitly as follows:

\begin{optprog}
{\bf (SP$_j$)} Minimize & \objective{\hat{f}_jt_j + C^V \sum_{h \in H_j}v_h - \sum_{h \in H_j}\sum_{k \in M}\sum_{i \in M} \hat{c}_{ki} x_{kjh}- \alpha_j} \nonumber\\
           & \sum_{k\in M} x_{ikh}-\sum_{k \in M} x_{kih} & = & 0 & \forall i \in M, h \in H_j, \label{sp-c-1}\\
           & y_{ikh} - L^V x_{ikh} & \leq & 0 & \forall i \in M, k \in M, h \in H_j, \label{sp-c-2}\\
           & \sum_{k \in M} y_{ikh} - \sum_{k \in M} y_{kih} + D_i\sum_{k \in M}x_{ikh} & = & 0 & \forall i \in I, h \in H_j, \label{sp-c-3}\\
           & \sum_{i \in M}\sum_{k \in M} T_{ik} x_{ikh} - L^T v_h & \leq & 0 & \forall h \in H_j, \label{sp-c-4}\\
           & x_{ikh} & \in & \{0,1\} & \forall i \in M, k \in M, h \in H_j, \label{sp-c-5}\\
           & y_{ikh} & \geq &  0 & \forall i \in M, k \in M, h \in H_j, \label{sp-c-6}\\           
           & t_{j}  & \in & \{0,1\}, & \mbox{and} \label{sp-c-7}\\
           & v_{h} & \in & \{0,1\} &  \forall h \in H_j, \label{sp-c-8}
\end{optprog}
where $M = I \cup \{j\}$ and 
\begin{eqnarray*}
\hat{f}_j & = & C^F_j-L^F_j \mu_j, \\
\hat{c}_{ki} & = & \left \{  \begin{array}{ll}
                 C^O T_{ki} - \pi_{i} + \mu_{j}D_{i} & \mbox{if $k \in M$, and $i \in I$}, \\
                 C^OT_{kj} & \mbox{if $k \in I$}.
                 \end{array}
              \right. \\
\end{eqnarray*}
The objective of each of the resulting single-facility
subproblems SP$_j$ is to generate sets of paths with least-cost assigned to 
{\em all vehicles} located at the facility, though without the constraint that 
customers appear exactly once. Because of this condition, we avoid calling
paths as routes for now.

% \begin{optprog}
%{\bf(SP_j)} Minimize & \objective{z=\hat{C}(\pi,\mu)U - \alpha} \label{obj-fullsp}\\
%subject to & \sum_{k\in I} X_{ikh}-\sum_{k \in I} X_{kih} & = & 0 & \forall i \in I, h \in H \nonumber\\
%           & Y_{ikh} - CapV \cdot X_{ikh} & \leq & 0 & \forall i,k \in I, h \in H\nonumber \\
%           & \sum_{k \in I} Y_{ikh} - \sum_{k \in I} Y_{kih} + D_i \sum_{k \in I}X_{ikh} & = & 0 & \forall i \in I, h \in H \nonumber\\
%           & \sum_{i \in I}\sum_{k \in I}TD_{ik}X_{ikh} - TL\cdot V_h  & \leq & 0 & h \in H \nonumber\\
%           & X_{ikh} & \in & \{0,1\} & \forall i,k \in I, h \in H\nonumber\\
%           & Y_{ikh} & \geq &  0 & \forall i,k \in I,  h \in H \nonumber \\
%           & V_h  & \in & \{0,1\} & \forall h \in H\nonumber\\
%           & T_j  & \in & \{0,1\} & \forall j \in J\nonumber
%\end{optprog}
%where $\hat{C}(\pi,\mu)$ is a vector of objective function coefficients 
%(constructed using (\ref{dw-red-cost-open})) for variable
%vector $U$ for given dual vectors $\pi$ and $\mu)$. For a
%given solution $\bar{U}=(\bar{X}^T,\bar{Y}^T, \bar{T}^T, \bar{V}^T)^T$, the objective function value is:
%\begin{eqnarray}
%z& = &\sum_{j \in J}(FC_j-CapF_j \mu_j) \bar{T}_{j}+ VF \sum_{h \in H}\bar{V}_h-\nonumber\\
%              & & \sum_{h \in H}\sum_{k \in I} (\sum_{i \in I} (OC TD_{ki} - \pi_{i} + \mu_{j}D_{i})\bar{X}_{kih} + OC\sum_{j \in J}TD_{ij}\bar{X}_{kjh}) - \alpha \nonumber
%\end{eqnarray}

To obtain integer feasible solutions to LRSP using MDW, we have to force
integrality restriction for the variable $\theta$. Thus, we add  
$\theta_q \in \{0, 1\}$ for all $q \in E$ to MDW. 

Now we prove that solving integer MDW is equal to solving SPP-LRSP. 
The similar forms of SPP-LRSP and integer MDW should now be evident, but to
rigorously show their equivalence, we need to dissect the relationship
between set $E_j$ (used to define MDW) and $P_j$ (used to define SPP-LRSP)
for a given facility $j \in J$.

%$P_j$ is the set of pairings (set of routes that include each customer
%at most once and can be served by one vehicle within the time limit) for 
%facility $j$, is defined in section \ref{set-part-form} for SPP-LRSP.

A member of set $E_j$, a feasible solution to SP$_j$, 
consists of a {\em collection} of sets of paths (starting at facility $j$
and ending at facility $j$) assigned to all vehicles located at 
facility $j$. It may be possible that some customers are visited 
multiple times by a single vehicle as well as different vehicles. 

As explained in Section \ref{set-part-form}, a member of set $P_j$, 
on the other hand, is a single set of routes (pairing) that can be 
assigned to any vehicle at facility $j$. A customer can only appear 
at most once in a pairing assigned to a vehicle.  

Some members of $E_j$ can be constructed by associating at most
$|H_j|$ members of set $P_{j}$ and some number of empty pairings (zero
vectors representing vehicles that are not used) to the vehicles
assigned to facility $j$. The members of $E_j$ that  include sets of 
paths assigned to a single vehicle and visiting some customers multiple times
can not be constructed using members of $P_{j}$. 

Now, by utilizing the integrality requirements from the original
problem and carefully eliminating the indices of symmetric solutions
from $E$, we get a much smaller set that we will show is in one-to-one
correspondence with collections of members of $P_j$ of cardinality at
most $|H_j|$. We proceed as follows:

% same customers are visited in a single route, or by vehicle..

\begin{enumerate}
\item First, to prove the necessity of adding 
  $\theta_q \in \{0, 1\}$ for all $q \in E$ to MDW, we 
  let $\hat{\theta}$ represent a solution to MDW for which the
  corresponding solution in the original space obtained by applying
  equations (\ref{XintermsTeta1})-(\ref{VintermsTeta1}) is feasible for the
  original problem G-LRSP. From (\ref{XintermsTeta1}) and
  (\ref{x-var}), along with the fact that $\mbox{x}^q_{ikh}$ is binary
  for $q \in E$, we get that $\mbox{x}^{q_1}= \mbox{x}^{q_2}$ 
  for all $q_1$ and $q_2$ such that $\hat{\theta}_{q_1} > 0$ 
  and  $\hat{\theta}_{q_2} > 0$ in the solution to MDW. 
  In other words, all  members $q$ of $E_j$ with $\hat{\theta}_q > 0$ 
  must correspond to set of paths visiting exactly the same customers in exactly 
  the same order. Similar results follow for $\mbox{t}^q$ and
  $\mbox{v}^q$ from (\ref{Tvartrans1}), (\ref{VintermsTeta1}), 
  and (\ref{t-var}), (\ref{dri-var}).
  Hence, we must in fact have $\hat{\theta}_q = 1$ for exactly
  one $q \in E_j$ for each $j \in J$ and so $\theta_q \in \{0, 1\}$
  for all $q \in E$.
\item Let $q \in E$ a member in which a vehicle visits at least one
  customer multiple times. Thus, $\mbox{x}_{ik_1h_1}^q = 1$ and
  $\mbox{x}_{ik_2h_1}^q = 1$ for some customer $i \in I$, 
  vehicle $h_1 \in H$ and nodes $k_1$, $k_2$ $\in N$ such that
  $k_1 \neq k_2$. This means that  $\sum_{k \in N}\mbox{x}_{ikh_1}^q > 1$
  and therefore  $b_{iq}> 1$. Because of constraint (\ref{spar1}) and  
  the binary restriction for $\theta_q$, a solution to integer MDW for which 
  the corresponding solution in the G-LRSP is feasible, $\theta_q$ must
  be 0. Therefore, we can remove such members from $E$. In order not 
  to generate these members, we can add the following constraint
  to SP$_j$.
  \begin{optprog}
    & \sum_{k \in M}x_{ikh} & \leq & 1 & \forall i \in N, h \in H_j \label{relax_sp1}
  \end{optprog}
  When we add constraint (\ref{relax_sp1}) to SP$_j$, the feasible solutions
  are restricted to sets of routes each of which assigned to a vehicle and 
  includes any customer at most once. Therefore, based on the definition 
  of a pairing in Section \ref{set-part-form},  we can say that the feasible
  solutions to augmented SP$_j$ are sets of pairings.
\item As the vehicles associated with a given facility $j$ are
  identical, a set of pairings can be assigned to vehicles in
  any arbitrary order. Hence, we obtain different members of $E_j$ that are 
  all equivalent from the standpoint of both feasibility and cost. 
  To eliminate superfluous equivalent members of $E_j$, we divide the 
  members of set $E_j$ into equivalence classes, where two members of $E_j$ 
  are considered equivalent if the set of customers assigned to the facility 
  and the partition of that set of customers defined by the pairings 
  (set of routes assigned to a vehicle) are identical. 
  It is clear that any two members of $E_j$ that are equivalent by this 
  definition will have exactly the same impact on both cost and feasibility. 
  
  In each equivalence class of $E_j$ there can be up to $|H_j|!$ members depending
  on the number of vehicles serving to non zero pairings in the solutions. 
  For example, let $q_1$ and $q_2$ be two members of $E_j$ and be in
  the same equivalence class. Then there exists a pair of vehicles $h_1$ and
  $h_2$ such that $\mbox{x}^{q_1}_{ikh_1} =\mbox{x}^{q_2}_{ikh_2}$, 
  $\mbox{x}^{q_1}_{ikh_2} =\mbox{x}^{q_2}_{ikh_1}$ $\forall i, k \in M$ and 
  at least one of $h_1$ and $h_2$ is serving a non zero pairing.
  
  %$q=(x_q^T, y_q^T, t_q^T, v_q^T)^T$
  %be a member of set $E$, and $h_1$ and $h_2$ $\in H_j$ be two vehicles located at facility $j$ for
  %some $j \in J$.
  %Assume that $v_{h_1q}=v_{h_2q}=1$ and $h_1$ and $h_2$ serve different non zero pairings.
  %Let $q'=(x^T_{q'}, y^T_{q'}, t^T_{q'}, v^T_{q'})^T$ be another solution
  %which is exactly same with $q$ except $x_{ikh_1q}=x_{ikh_2q'}$, $x_{ikh_2q}=x_{ikh_1q'}$
  %and $y_{ikh_1q}=y_{ikh_2q'}$, $y_{ikh_2q}=y_{ikh_1q'}$ $\forall i, k \in I$.
  %%and $\bar{x}_{ikh_1q}>0$ and $\bar{x}_{sth_2q}>0$ for some $i, k, s, t \in I$.
  %It is clear that $q$ and $q'$ are in same symmetry group. 
  %Under the assumption of identical vehicles, it is indifferent to choose solution $q$ 
  %or solution $q'$. Therefore, we keep only one member of a symmetry group and eliminate
  %the others from set $E$. This process  reduces the degeneracy in
  %master problem.       

  We form equivalence classes from which all but one member may safely be 
  eliminated from $E_j$.
  %%%%%%%%%%%% from lrp paper %%%%%%%%%%%%%%%%%%%%%%%%%%%%%%%%%%%%%%%%%%%%%
  %\item First, as the vehicles associated with a given facility $j$ are
  %  identical, a set of pairings from $P_j$ can be assigned to vehicles in
  %  any arbitrary order and each route can also visit the customers in
  %  either clockwise or counterclockwise order. Hence, we obtain
  %  different members of $E_j$ that are all equivalent from the standpoint
  %  of both feasibility and cost. To eliminate superfluous equivalent 
  %  members of $E_j$, we divide the members of set $E_j$ into equivalence 
  %  classes, where two members of $E_j$ are considered equivalent if 
  %  the set of customers assigned to the facility and the partition of that 
  %  set of customers defined by the pairings (set of routes assigned to
  %  a vehicle) are identical. It is clear that any two members of $E_j$ that
  %  are equivalent by this definition will have exactly the same impact on
  %  both cost and feasibility. We then form equivalence classes from which
  %  all but one member may safely be eliminated from $E_j$.
\item 
  %It then follows easily that 
  Index $q$ can be removed from $E_j$
  if there exist vehicles $h_1, h_2 \in H_j$ such that
  $\mbox{x}_{ikh_1}^q = \mbox{x}_{ikh_2}^q$ $\forall i, k \in M$, where
  $\mbox{x}_{ikh_1}^q>0$ for some $i, k \in M$ (i.e., this is not an
  empty route). In this case, vehicles $h_1$ and $h_2$ define exactly
  the same sets of routes, which means that $b_{iq}> 1$ for some $i \in
  I$. Because of constraint (\ref{spar1}), such a solution must have
  $\hat{\theta}_q=0$.
\item From (\ref{Tvartrans1}) and (\ref{rel-2a}), we can
  conclude that if $\hat{\theta}_q =1$ and we have $\mbox{x}_{ikh}^q>0$
  for some $i, k \in M$ and $h \in H$, then $\mbox{t}^q$ must be $1$.
  Hence, we can eliminate any solution for which $\mbox{t}^q = 0$
  that does not correspond to a zero solution (i.e., closed facility).
  All of this allows us to rewrite (\ref{rel-2a}) in the form
  \begin{equation}
     -\sum_{q \in E_j} \sum_{i \in I} b_{iq} D_i \theta_q  
     + L^F_j \sum_{q \in E_j} \theta_q \geq 0 \ \ \forall j \in J. \label{rel-2-new-a}
  \end{equation}
\item A member $q \in E_j$ %$=(x_q^T, y_q^T, t_q^T, v_q^T)^T$ 
  can be removed if there exists a vehicle $h \in H_j$
  such that  $\mbox{v}^q_{h} = 1$ and  $\mbox{x}^q_{ikh} = 0$ 
  for all $i,k \in M$ (meaning that zero vector is assigned to vehicle $h$
  even though we pay the fixed cost of that vehicle).
  In an optimal solution to integer MDW, $\theta_{q'} \geq \theta_q = 0$
  where $q, q' \in E_j$, $\mbox{x}^q = \mbox{x}^{q'}$, $\mbox{y}^q = \mbox{y}^{q'}$,
  $\mbox{t}^q = \mbox{t}^{q'}$, and $\mbox{v}^q = \mbox{v}^{q'}$
  except that $\mbox{v}^{q'}_{h} = 0$ since $\tilde{C_{q}} \geq \tilde{C_{q'}} $.
  %from set $E$ if there exists some set 
  %$\mathcal{F_H}=\{h \in H| v_{hq}=1, x_{ikhq}=0 \ \  \forall i,k \in I\}$. 
  %This means that we do not assign vehicles in set  $\mathcal{F_H}$ to any
  %pairings (i.e., zero vector is assigned), but we pay the fixed cost of vehicles. A member similar to $q$
  %except $v_{hq}=0$ $\forall h \in \mathcal{F_H}$ is always preferred by the master problem to $q$
  %because of the cost, therefore we can remove members satisfying this condition. 
\item Finally, we can remove index $q$ corresponding zero column (i.e., 
  $\mbox{x}^q = 0$, $\mbox{y}^q = 0$, $\mbox{v}^q = 0$, and $\mbox{t}^q = 0$)
  from $E_j$. Thus, we can update constraint (\ref{conv2a}):
  \begin{optprog}
    & \sum_{q \in E_j} \theta_q  & \leq & 1 & \forall j \in J,\label{conv2b} 
  \end{optprog}
\end{enumerate}  

%%%%%%%%%%%%%%%%%%%%%%%%%%%%%% COMMENT %%%%%%%%%%%%%%%%%%%%%%%%%%%%%%%%%%%%%%%%%%%%%%%%%%%%%%%%%%%%%
\begin{comment}
It is clear that (Sub-Prob) decomposes into $|H|$ independent subproblems, 
one for each vehicle.
Based on our assumption that the vehicles are identical, the subproblem for 
each vehicle located at the same facility 
is identical. Therefore, by solving $|J|$ smaller subproblems, one for each facility, 
we can enumerate the integer feasible points for (Sub-Prob). The decomposed subproblem
(Dec-Sub-Prob) for each facility $j \in J$ is formulated as follows: 
 \begin{optprog}
{\bf(Dec-Sub-Prob)} Minimize & \objective{\sum_{i \in I} \sum_{k \in I} \hat{c}_{ik} X_{ik}+K }\label{de-subprob-obj}\\
subject to & \sum_{k\in I} X_{ik}-\sum_{k \in I} X_{ki} & = & 0 & \forall i \in I\nonumber\\
           & Y_{ik} - CapV \cdot X_{ik} & \leq & 0 & \forall i,k \in I\nonumber \\
           & \sum_{k \in I} Y_{ik} - \sum_{k \in I} Y_{ki} + D_i \sum_{k \in I}X_{ik} & = & 0 & \forall i \in N \nonumber\\
           & \sum_{i \in I}\sum_{k \in I}t_{ik}X_{ik}  & \leq & TL \nonumber\\
           & X_{ik} & \in & \{0,1\} & \forall i,k \in I \nonumber\\
           & Y_{ik} & \geq &  0 & \forall i,k \in I\nonumber 
\end{optprog}
where $\hat{c}$ is a vector for the reduced cost coefficients of variable $X$, $K$ is a constant, 
including costs corresponds to one vehicle and facility $j$
and $I=N\cup \{j\}$. $\hat{c}$ and $K$ are defined using (\ref{dw-red-cost-open}):
\begin{eqnarray*}
\hat{c}_{ki} & = & \left \{  \begin{array}{ll}
                 (OC t_{ki} - \pi_{i} + \mu_{j}D_{i})  & \mbox{if $k \in I$, and $i \in N$} \\
                 (OC t_{ki}) & \mbox{if $k \in I$ and $i=j$}
                         \end{array}
              \right. \\
K & = & (FC_j-CapF_j\mu_j)+ VF - \alpha 
\end{eqnarray*}
 
(Dec-Sub-Prob) for facility $j \in J$ generates feasible pairings 
%(set of vehicle routes that can be seved by one vehicle within the time limit) 
that is represented with notation $P_j$ $\forall j \in J$ in section (\ref{set-part-form}). 
(Dec-Sub-Prob) is an elementary shortest path problem with time and capacity constraints.
Since (Dec-Sub-Prob) is 
obtained by decomposing (Sub-Prob) into independent problems, there is a relationship between
sets  $E$ and $P_j$  $\forall j \in J$. 
%We can construct members of set $E$ using the members of sets $P_j$ $\forall j \in J$.
A member of set $E$ is a feasible solution to (Sub-Prob) 
that  presents sets of open and closed facilities, 
sets of vehicles used and a set of routes (pairings) for each vehicle  
in that solution. A member of set $P_j$ is a feasible solution to (Dec-Sub-Prob) that presents a pairing,
a set of routes.  
Therefore, in a member of $E$, pairings assigned to vehicles located at facility $j$ are members of set 
$P_{j}$ $\forall j \in J$.  In addition to this relationship, 
picking a member of set $P_j$ implies that vehicle $h \in H_j$ that is 
assigned to the pairing is used, $\forall j \in J$. 
The zero vector is feasible for (Dec-Sub-Prob), but we assume that sets $P_{j}$
$\forall j \in J$ do not include zero vector. We consider assignment of zero
vector to vehicles separately.
%We simply say that if a vehicle is not used in a solution, 
%then the  pairing corresponds to zero vector is assigned to that vehicle. 

Same subsets of pairings $P_j$, $\forall j \in J$ can be used in many members of set $E$. 
To make this relationship one to one, we can do preprocessing and restrict set $E$.
We remove some members of $E$ whose corresponding variables must be set to zero 
in an optimal solution of (Master-DW-1). We can state following preprocessing 
rules to restrict set $E$:

\begin{enumerate}
\item A member $q=(x_q^T, y_q^T, t_q^T, v_q^T)^T$ that is a solution to (Sub-Prob) 
can be removed from set $E$ if there exists $h \in H_j$ and $j \in J$ such that $v_{hq}=1$ and 
$t_{jq}=0$. This
 means that we use vehicle $h$ located at facility $j$, but the facility is closed. The relationship between
variables $X$, $Y$ and $V$ is forced in (Sub-Prob). If $v_{hq}=1$ then there exists some $i \in N$ such that 
$x_{jihq}=1$ and $y_{jihq}>0$. Based on master problem constraints (\ref{rel-2}),
if $y_{jihq}>0$ for some $i \in N$ and $t_{jq}=0$, then the master problem variable $\theta_q$
corresponds to member $q=(x_q^T, y_q^T, t_q^T, v_q^T)^T$ must be zero. 
Therefore, we can remove members of $E$ satisfying this condition.

%\item A member $q=(\bar{x}_q^T, \bar{y}_q^T, \bar{t}_q^T, \bar{v}_q^T)^T$ can be removed 
%from set $E$ if there exists some set 
%$\mathcal{F_M}=\{j \in M| \bar{t}_{jq}=1, \bar{v}_{hq}=0 \ \ \forall h \in H_j\}$. 
%This means that we do not use any vehicle located at facility $j \in \mathcal{F_M}$, but the facility is open. 
%There exists a member $q'=(\bar{x}^T_{q'}, \bar{y}^T_{q'}, \bar{t}^T_{q'}, \bar{v}^T_{q'})^T$
%in $E$ such that $\bar{x}_q=\bar{x}_{q'}$, $\bar{y}_q=\bar{y}_{q'}$, $\bar{v}_q=\bar{v}_{q'}$, and 
%$\bar{t}_{jq}=\bar{t}_{jq'}$ $\forall j \in M \backslash \mathcal{F_M}$ and $\bar{t}_{jq'}=0$ 
%$\forall j \in \mathcal{F_M}$. 
%Let $\bar{c}_q$ and $\bar{c}_{q'}$ be the costs of these solutions. We know that 
%$\bar{c}_{q'}=\bar{c}_{q}-\sum_{j \in \mathcal{F_M}}FC_j$. 
%Since we want to find a solution with minimum cost, the master problem 
%prefers solution $q$ to $q'$. Therefore, variable $\theta_q$ 
%corresponds to member $q=(\bar{x}_q^T, \bar{y}_q^T, \bar{t}_q^T, \bar{v}_q^T)^T$ can be set to zero. 
%We can remove members of $E$ satisfying this condition.

\item A member $q=(x_q^T, y_q^T, t_q^T, v_q^T)^T$ can be removed 
from set $E$ if there exists some set 
$\mathcal{F_H}=\{h \in H| v_{hq}=1, x_{ikhq}=0 \ \  \forall i,k \in I\}$. 
This means that we do not assign vehicles in set  $\mathcal{F_H}$ to any
pairings (i.e., zero vector is assigned), but we pay the fixed cost of vehicles. A member similar to $q$
except $v_{hq}=0$ $\forall h \in \mathcal{F_H}$ is always preferred by the master problem to $q$
because of the cost, therefore we can remove members satisfying this condition. 
\item A member $q=(x_q^T, y_q^T, t_q^T, v_q^T)^T$ can be removed 
from set $E$ if there exists two vehicles located at the same facility
$h_1, h_2 \in H_j$, for some $j \in J$ such that $x_{ikh_1q}=x_{ikh_2q}$ 
$\forall i, k \in I$ and  $x_{ikh_1q}>0$ for some $i, k \in I$. This means that 
vehicles $h_1$ and $h_2$ serve exactly the
same set of routes which leads to $\sum_{h \in H_j}\sum_{k \in I}x_{ikhq}=g_{jiq}>1$.
Because of constraint (\ref{spar}), such a solution can not be selected by master problem,
i.e., $\theta_q=0$.
%a customer $i \in N$ and a facility $j \in J$ such that 
%$\sum_{h \in H_j}\sum_{k \in I}\bar{x}_{ikh}=g_{jiq}>1$. This means that customer
%$i$ is served more than one times from facility $j$. Such a solution cannot be selected
%by master problem, i.e., $\theta_q=0$.
\item Since vehicles located at the same facility are identical, assignment of same set
of pairings to different set of vehicles leads to different solutions that are symmetric
and have same cost. We can divide the members of set $E$ into symmetry groups. 
In each symmetry group there can be up to $|H_j|!$ members depending
on the number of vehicles serving to non zero pairings in the solutions. Let 
$q=(x_q^T, y_q^T, t_q^T, v_q^T)^T$
be a member of set $E$, and $h_1$ and $h_2$ $\in H_j$ be two vehicles located at facility $j$ for
some $j \in J$.
Assume that $v_{h_1q}=v_{h_2q}=1$ and $h_1$ and $h_2$ serve different non zero pairings.
Let $q'=(x^T_{q'}, y^T_{q'}, t^T_{q'}, v^T_{q'})^T$ be another solution
which is exactly same with $q$ except $x_{ikh_1q}=x_{ikh_2q'}$, $x_{ikh_2q}=x_{ikh_1q'}$
and $y_{ikh_1q}=y_{ikh_2q'}$, $y_{ikh_2q}=y_{ikh_1q'}$ $\forall i, k \in I$.
%and $\bar{x}_{ikh_1q}>0$ and $\bar{x}_{sth_2q}>0$ for some $i, k, s, t \in I$.
It is clear that $q$ and $q'$ are in same symmetry group. 
Under the assumption of identical vehicles, it is indifferent to choose solution $q$ 
or solution $q'$. Therefore, we keep only one member of a symmetry group and eliminate
the others from set $E$. This process  reduces the degeneracy in
master problem.       
\end{enumerate}   
\end{comment} 
%%%%%%%%%%%%%%%%%%%%%%%%%%%%%%%%%%%%%%%%%%%%%%%%%%%%%%%%%%%%%%%%%%%%%%%%%%%%%

If we restrict set $E_j$ according to the rules described above and call the
restricted set $\bar{E}_j$ for each $j \in J$, then we can finally
conclude the following.
\begin{prop}
\label{dw-prop-p1}
There is a one-to-one correspondence between nonempty subsets of $P_j$ with
cardinality less than $|H_j|$ and members of $\bar{E}_j$.
\end{prop}
\begin{proof}
  The proof follows easily from the definition of sets $P_j$ and $\bar{E}_j$,
  and the restriction rules. 

%  \zaupdate{Do I need to include the longer proof}
  \begin{comment}
    Members of set $\bar{E}_j$ are obtained using SP$_j$ augmented with
    constraint set (\ref{relax_sp1}). All constraints 
    of this system except (\ref{sp-c-7}) (binary restriction for facility
    location variable) are indexed by the set $H_j$ and it is therefore 
    immediate that this system decomposes by vehicle. Based on the assumption
    that all vehicles are identical, we can drop the index h
    For each vehicle  $h \ in H_j$,
    %  candidate facility $j \in J$, the set of integer solutions of the 
    %  decomposed subproblem can be represented by the set of vectors defined as
    %  \begin{eqnarray}
    %    & & \{ \ \  (x, t, v)^j \in \Z^{(|N|\times |N| \times |H_j|) \times \{0, 1\} \times(|H_j|)}, 
    %    (y)^j \in \R^{(|N|\times |N| \times |H_j|)}  \mid  \ \  (x, y, t, v)^j \mbox{ satisfies } \nonumber\\
    %    & &  \ \ (\ref{enterleave})_j, 
    %    (\ref{VehCap})_j, (\ref{flow-conv})_j, (\ref{timeLim})_j, (\ref{fix-0})_j, 
    %    (\ref{x-var})_j, (\ref{y-var})_j, (\ref{t-var})_j, (\ref{dri-var})_j \}, \label{subsys} 
    %  \end{eqnarray}
    %which is described only by those constraints associated with facility
    %$j$ and specifies values only for those variables associated with
    %vehicles assigned to facility $j$. Hence, the index $j$ for the
    %constraints represents the set of constraints associated with facility
    %$j$. 
  \end{comment}
  \qed 
\end{proof}

%%%%%%%%%%%%%%%%%%%%%%%%%%%%%%%%%%%%%%%%%%%%%%%%%%%%%%%%%%%%%%%%%%%%%%%%%%%%%
\begin{comment}
\begin{prop}
\label{dw-prop-p1}
Let $\bar{P}_j$ be a subset of $P_j$ such that $n_j=|\bar{P}_j|\leq |H_j|$ $\forall j \in J$
and let $\mathcal{P}=\cup_{j \in J}\bar{P}_j$.
Let $q=(x_{q}^T, y_{q}^T, t_{q}^T, v_{q}^T)^T$ be a member of set $E_r$. 
There is a one to one correspondence
between every possible $\mathcal{P}$ and vector  $(x_{q}^T, y_{q}^T, v_{q}^T)^T$ 
$\forall q \in E_r$.
\end{prop}
\begin{proof}
Based on the preprocessing rule 4, $E_r$ has one member representing each symmetry group
in $E$. We can set a specific rule to determine such a representative member for $E_r$. 
We can assume that pairings in sets $P_j$ $\forall j \in J$, and vehicles in sets 
$H_j$ $\forall j \in J$ are ordered. We only keep $q \in E_r$ that has the following structure. 
For a given $\bar{P}_j$, let $p_{i_1}, p_{i_2}, .., p_{i_{n_j}} \in \bar{P}_j$ and 
$h_1, h_2, h_3, .., h_{n_j} \in H_j$
such that the indices of pairings satisfy the relation $i_1 < i_2 < .. < i_{n_j}$
$\forall j \in J$.
We can simply construct one member of symmetry group corresponds to $\mathcal{P}$ by assigning
$p_{i_1}$ to $h_1$, $p_{i_2}$ to $h_1$, .., $p_{i_{n_j}}$ to $h_{n_j}$ $\forall j \in J$
and assigning zero vector to the remaining vehicles. 

Let $p$ be a pairing in set $P_j$. $p$ includes a set of arcs that form a set of routes.
Let $\mathcal{I}_p$ be the set of arcs included in pairing $p$.
In a solution $q \in E_r$, assigning $p \in P_j$ to a vehicle $h \in H_j$ for $j \in J$
means that $x_{ikhq}=1$ $\forall (i,k) \in \mathcal{I}_p$ and $x_{ikhq}=0$ 
$\forall (i,k) \in I \backslash \mathcal{I}_p$. We can easily construct $y_q$ vector
using $x_q$ vector and demand information. $y_q$ vector keeps the vehicle load necessary 
to serve the remaining customers on a route. Therefore, based on preprocessing
rule 4, set  $\mathcal{P}$ corresponds to a unique vector  
$(x_{q}^T, y_{q}^T)^T$.

From preprocessing rule 2, $v_{hq}=1$ $\Leftrightarrow$  $x_{ikhq}=1$ 
for some $(i,k) \in I$, $\forall h \in H$ and $\forall q \in E_r$. Therefore, set 
$\mathcal{P}$ leads to vector $v$ such that $v_{h_iq}=1$ for 
$i=1,..,n_j$ and $v_{h_iq}=0$ for $i=n_j+1, ..,|H_j|$ $\forall j \in J$. 

Similarly, we can construct members of set $\mathcal{P}$ from a given
vector $(x_{q}^T, y_{q}^T)^T$ for some $q \in E_r$ by identifying pairings
corresponds to set of arcs visited by each vehicle. Preprocessing rules
3 and 4 guarantee that there is not any other vector $(x_{q'}^T, y_{q'}^T)^T$
for some $q' \in E_r$ corresponds to same set $\mathcal{P}$.
 \qed 
\end{proof}
%\begin{enumerate}
%\item [C1] $\bar{t}_{jq}=1$ $\Leftrightarrow$  $\bar{v}_{hq}=1$ for some $h \in H_j$, $\forall j \in M$. 
%\item [C2] $\bar{t}_{jq}=0$ $\Leftrightarrow$  $\bar{v}_{hq}=0$ for all  $h \in H_j$, $\forall j \in M$. 
%\item [C3] $\bar{v}_{hq}=1$ $\Leftrightarrow$  $\bar{x}_{ikhq}=1$ for some $(i,k) \in I$.
%\item [C4] $p \in P_j$ $\forall j \in M$ can be assigned to at most one vehicle.
%\item [C5] $q$ is the only member in $E_r$ from the symmetry group that $q$ belongs in $E$. 
%We can assume that pairings in sets $P_j$ $\forall j \in M$, and vehicles in sets 
%$H_j$ $\forall j \in M$ are ordered. We use this assumption to delete members in a same symmetry 
%group in $E$. We only keep $q \in E_r$ that has the following structure. 
%Let $p_{i_1}, p_{i_2}, .., p_{i_n} \in P_j$ and $h_1, h_2, h_3, .., h_n \in H_j$ for some
%$n \leq |H_j|$ and $j \in M$ such that the indices of pairings satisfy the relation $i_1 < i_2 < .. < i_n$.
%We can simply construct one member of symmetry group by assigning
%$p_{i_1}$ to $h_1$, $p_{i_2}$ to $h_1$, .., $p_{i_n}$ to $h_n$ and assigning zero vector to the 
%remaining vehicles which means that vehicles $h_{n+1}, .., h_{|H_j|}$ are not used.  
%\end{enumerate}   
We form master problem (Master-DW-1') which has the same set of constraints with (Master-DW-1)
and uses set $E_r$ instead of set $E$. Now we can state the following relationship between
model (Master-DW-1') and (SPP-LRS) defined in section \ref{part1:sp_form}.
%It is true that we can add this trivial upper bound for the number of vehicles for
%each facility to our SPP-LRS formulation defined in section \ref{part1:sp_form}. 
%Therefore, Master-DW-2 is equal to our SPP-LRS formulation plus the upper bound for
%the total number of vehicles used for each facility.  
\begin{prop}
\label{dw-prop-p2}
Every integer feasible point for (SPP-LRS)  can be one to one mapped to an 
integer feasible point for (Master-DW-1') with same objective function value. 
\end{prop}
\end{comment} 
%%%%%%%%%%%%%%%%%%%%%%%%%%%%%%%%%%%%%%%%%%%%%%%%%%%%%%%%%%%%%%%%%%%%%%%%%%%%%

By replacing set $E_j$ with $\bar{E}_j$
for all $j \in J$ in MDW, as well as replacing (\ref{rel-2a}) with
(\ref{rel-2-new-a}) and (\ref{conv2a}) with (\ref{conv2b}),
and adding the constraint $\theta_q \in \{0, 1\}$ for
all $q \in E$, we obtain a new (equivalent) formulation MDW'. Then,
we finally have the equivalence of SPP-LRSP and integer MDW as follows.
\begin{prop}
\label{dw-prop-p2}
There is a one-to-one correspondence between solutions to SPP-LRSP
and solutions to MDW' such that corresponding solutions also have the
same objective function value. Thus, SPP-LRSP, MDW', and integer MDW 
are all equivalent.
\end{prop}
\begin{proof}
  First, we construct a one-to-one  mapping to establish the 
  correspondence between solution sets of SPP-LRSP and MDW'.
  We define mapping $\Psi:(t, z) \rightarrow \theta$.
  Let $(\bar{t}, \bar{z})$ be a feasible solution to SPP-LRSP
  where $\bar{t}_j = 1$ if facility $j$ is open, $\bar{z}_p = 1$ 
  if pairing $p$ is selected and assigned to a vehicle 
  for all $p \in P_j$ and $j \in J$, 
  then $\bar{\theta}=\Psi(\bar{t}, \bar{z})$ be the corresponding
  solution in MDW' where $\bar{\theta}_q = 1$ if 
  solution $q$  is selected for all $q \in \bar{E}_j$, $j \in J$.
  
  Since we can decompose variables in both formulations by facility,
  we can decompose mapping $\Psi$ by facility and define
  mapping $\Psi_j$ indexed by facility. 
  $\Psi = (\Psi_j)_{j \in J}$ and 
  $\bar{\theta}_{\bar{E}_j} = \Psi_j (\bar{t}_j, \bar{z}_{P_j})$
  where notation $x_S$ represents a vector composed of
  the components of vector $x$ indexed with set $S$. 
  %construct a mapping indexed by facility $j$ $\forall j \in J$.
  Let $\bar{J}=\{j \in J | \bar{t}_j=1\}$, 
  $\bar{P}_j=\{p \in P_j | \bar{z_p}=1\}$ 
  $\forall j \in \bar{J}$ and 
  $\bar{P}_j=\emptyset$ $\forall j \in J \backslash \bar{J}$ (because
  of constraint (\ref{spp_con2})). 
  For some facility $j \in J$, we have three types of solutions, which
  changes the definition of 
  $\bar{\theta}_{\bar{E}_j} = \Psi_j(\bar{t}_j, \bar{z}_{P_j})$. %conditions:
  \begin{itemize}
  \item [i.] Zero solution, $\bar{t}_j = 0$ and $\bar{P}_j=\emptyset$, thus
    $\bar{z}_p = 0$ for all $p \in P_j$. Then, 
    $\Psi_j(\bar{t}, \bar{z}_{P_j}) =$ $\bar{\theta}_{\bar{E}_j} = 0$,
    i.e., $\bar{\theta}_q = 0$ for all $q \in \bar{E}_j$ 
    (since we removed zero solution from $\bar{E}_j$ in reduction
    rule 7).
  \item [ii.] Nonzero but empty solution, $\bar{t}_j = 1$ and 
    $\bar{P}_j=\emptyset$. Then, $\bar{\theta}_1 = 1$ such that
    $\mbox{x}^1 = 0$, $\mbox{y}^1 = 0$, $\mbox{v}^1 = 0$
    and $\mbox{t}^1 = \bar{t}_j = 1$, and $\bar{\theta}_q = 0$
    for all $q \in \bar{E}_j$ and $q \neq 1$.
  \item [iii.] $\bar{t}_j = 1$ and  $\bar{P}_j \neq \emptyset$.
    For each facility we can choose at most $|I| = |H_j|$ pairings,
    i.e., $\sum_{p \in P_j}\bar{z}_p = |\bar{P}_j| \leq |H_j|$.
    Based on Proposition \ref{dw-prop-p1}, there is one member 
    $k \in \bar{E}_j$ 
    (in which $|\bar{P}_j|$ vehicles assigned one pairing
    from set $\bar{P}_j$, $|H_j|-|\bar{P}_j|$ vehicles
    assigned zero solution, and $\mbox{t}^k_j = 1$) corresponds to 
    subset $\bar{P}_j$. Thus, the corresponding solution in MDW' is 
    $\bar{\theta}_k = 1$ and $\bar{\theta}_q = 0$ for all 
    $q \in \bar{E}_j$ and $q \neq k$. We use notation $k \Tr \bar{P}_j$
    to define this one-to-one correspondence between member $k$ and
    pairing set $\bar{P}_j$ defined in Proposition \ref{dw-prop-p1}.
  \end{itemize} 
  It follows that  $\Psi_j$ is one-to-one from construction, therefore
  $\Psi$ is also  one-to-one. It is clear that we can easily construct 
  the one-to-one reverse mapping,  
  $(\bar{t}_j, \bar{z}_{P_j})=\Psi_j^{-1} (\bar{\theta}_{\bar{E}_j})$ in 
  a similar way. 

  From the construction of $\Psi_j$ for each solution type, we can 
  write the following relations between two formulations' variables 
  and parameters.
  \begin{optprog}
  %\begin{eqnarray*}
    & t_j & = & \sum_{q \in \bar{E}_j}\theta_q  & \forall j \in J, \label{rel-a-b1}\\
    & b_{iq} & = & \sum_{p \in S: q \Tr S} a_{ip} & \forall i \in I, \forall q \in \bar{E}_j, \label{rel-a-b2}\\
    & \tilde{C_{q}} & = & C^F_j + \sum_{p \in S:q \Tr S } C_p & \forall q \in \bar{E}_j, \forall j \in J  \label{rel-a-b3}                           
    \end{optprog} 
  %\end{eqnarray*}

  If $(\bar{t}, \bar{z})$ is feasible, then 
  $\bar{\theta}=\Psi(\bar{t}, \bar{z})$ is also feasible. This can be seen 
  from the similarity between the structures of SPP-LRSP and MDW' and
  from (\ref{rel-a-b1}) and (\ref{rel-a-b2}). Finally, the
  equivalence of objective function values follows from
  (\ref{rel-a-b3}).

  %Let $z_{SPP}$ be the objective function
  %value of SPP-LRS for solution  $(\bar{t}, \bar{z})$ and $z_{DW}$ be the objective function
  %value of MDW' for solution  $\bar{\theta}=\Psi(\bar{t}, \bar{z})$. 
  %Let $J_i$, $J_{ii}, J_{iii}$ be a partition of facility
  %set including facilities in solution types (i), (ii), and (iii). 
  %The objective function coefficients of facilities in set $J_i$
  %are zero in   
  %solution types (i) and (ii), it is obvious that the contributions of 
  %objective function values for facility $j$ are equal.    
  %\begin{optprog}
  % & z_{SPP} & = & \sum_{j \in J} C^F_j \bar{t}_{j} + \sum_{j \in J} \sum_{p\in P_j} C_p \bar{z}_p \nonumber \\
  %  &        & = & \sum_{j \in J} (C^F_j \bar{t}_{j} + \sum_{p\in \bar{P}_j} C_p)  \nonumber %\\
  %\end{optprog}
\qed
  \end{proof}

Proposition \ref{dw-prop-p2} proved that applying Dantzig Wolfe decomposition
to G-LRSP,  is equal to solving SPP-LRSP.


%%%%%%%%%%%%%%%%%%%%%%%%%%%%%%%%%%%%%%%%%%%%%%%%%%%%%%%%%%%%%%%%%%%%%%%%%%%%%
\begin{comment}
\begin{proof}
We represent this mapping using notation $\Psi:(T,Z) \rightarrow \theta$. 
Let $(\bar{T}, \bar{Z})$ be a feasible solution to (SPP-LRS) 
and $\bar{\theta}=\Psi(\bar{T}, \bar{Z})$.
Based on constraint (\ref{conv2}) in (Master-DW-1)), a feasible solution
$\bar{\theta}$ must satisfy the following condition: there is one
$\bar{q}\in E_r$ such that $\bar{\theta}_{\bar{q}}=1$ and $\bar{\theta}_q=0$ 
$\forall q \in E_r \backslash \{\bar{q}\}$. Therefore,  to prove the 
proposition, it is enough to prove that 
$(\bar{T}, \bar{Z})$ can be one to one mapped to a member $\bar{q} \in E_r$.

In (SPP-LRS), $Z_p$ is a binary variable showing that pairing $p \in P_j$ $\forall j \in M$ 
is selected or not. A feasible solution $\bar{Z}$ defines a set of pairings selected from
each set $P_j$ $\forall j \in J$.   A pairing can only be selected
at most once, and at most $|H_j|$ pairings are selected from each facility $j \in J$
(constraint (\ref{vehicle-bound})).
Let $\bar{J}=\{j \in J | \bar{T}_j=1\}$, $\bar{P}_j=\{p \in P_j | \bar{Z_p}=1\}$ 
$\forall j \in \bar{J}$ and $\bar{P}_j=\emptyset$ $\forall j \in J \backslash \bar{J}$.
Let $\mathcal{P}=\cup_{j \in J}\bar{P}_j$ and  
$\bar{q}=(x_{\bar{q}}^T, y_{\bar{q}}^T, t_{\bar{q}}^T, v_{\bar{q}}^T)^T$.

From proposition \ref{dw-prop-p1}, $\mathcal{P}$ can be one to one mapped to
vector $(x_{\bar{q}}^T, y_{\bar{q}}^T, v_{\bar{q}}^T)^T$ that is included in 
some $\bar{q} \in E_r$. Then, map $\bar{T}$ to $t_{\bar{q}}$ using one to one
mapping: 
$t_{j\bar{q}}=\bar{T}_j$, i.e., $t_{j\bar{q}}=1$ $\forall j \in \bar{J}$ and  
$t_{j\bar{q}}=0$ $\forall j \in J \backslash \bar{J}$. This proves that the mapping
$\Psi$ is one to one. 
%\begin{itemize}
%\item [S1] 
%$\bar{t}_{j\bar{q}}=\bar{T}_j$, i.e., $\bar{t}_{jq}=1$ $\forall j \in \bar{M}$ and  
%$\bar{t}_{jq}=0$ $\forall j \in M \backslash \bar{M}$.
%\item [S2] Let $1, 2, .., |H_j|$ be the indices of vehicles located at facility $j$,
%$\forall j \in M$. $\bar{v}_{i}=1$ $\forall i=1,..,|\bar{P}_j|$ and $\bar{v}_{i}=0$ $\forall i=|\bar{P}_j|+1, .., |H_j|$ 
%$\forall j \in M$.  
%\item [S3] Let $p_{i_1}, p_{i_2}, .., p_{i_{|\bar{P}_j|}} $ be the ordering of
%pairings in $\bar{P}_j$ $\forall j \in \bar{M}$.  
%We assign $p_{i_1}$ to vehicle 1, $p_{i_2}$ to 
%vehicle 2, .., $p_{i_{|\bar{P}_j|}}$ to vehicle $|\bar{P}_j|$ located at facility $j$, 
%$\forall j \in \bar{M}$ and we assign zero vector to the remaining vehicles.
%\end{itemize} 

Now, we show that constructed $\bar{q}$ is a member of $E_r$. 
A feasible solution $(\bar{T}, \bar{Z})$ satisfy that  if $\bar{Z}_p=1$ for some $p \in P_j$,
then $\bar{T}_j$ must be 1 $\forall j \in J$. $\bar{T}_j$ can be 0 or 1 if   $\bar{Z}_p=0$   
$\forall p \in P_j$ (i.e.,$\bar{P}_j=\emptyset$). There cannot be a case where $\bar{Z}_p=1$ 
for some $p \in P_j$ and $\bar{T}_j=0$ for some $j \in J$. 
If $\bar{Z}_p=1$ for some $p \in P_j$,
then $|\bar{P}_j|\geq 1$, therefore mapping in  proposition \ref{dw-prop-p1} sets 
$v_{h\bar{q}}=1$ for some $h \in H_j$, $j \in J$. Preprocessing
rule 1 guarantees that there is not any member of $E_r$ satisfying 
$t_{j\bar{q}}=0$ when $v_{h\bar{q}}=1$ for some $h \in H_j$, $j \in J$.
Therefore, mapping from  solution $(\bar{T}, \bar{Z})$ leads
to a member of $E_r$. 

%and  $\forall j \in J$. Construction rules S1 and S2 define vectors $t_q$ and $v_q$, and any
%relation between these vectors is feasible with respect to (Sub-Prob) and preprocessing
%rule 1 eliminates solutions satisfying $t_{qj}=0$ and $v_{qh}=1$ for some $h \in H_j$ and $j \in J$
%(which would be created if  $T_j=0$ and $Z_p=1$ for some $p \in P_j$, $j \in J$). Therefore, 
%$q$ constructed using rules S1 and S2 is in set $E_r$ and there is not any member of
%$E_r$ including vectors $t$ and $v$ that cannot be constructed using  $(\bar{T}, \bar{Z})$.

%Every pairing (set of routes) $p \in P_j$ $\forall j \in J$ (generated using (Dec-Sub-Prob)) 
%is also feasible 
%for system (Sub-Prob) because of the relationship between models (Sub-Prob) and (Dec-Sub-Prob).
%Therefore, rule S3 leads $q$ which is a member of $E$. In addition, preprocessing
%rule 4 only keeps one specific member from each symmetry group, rule S3 lead $q$ which is a member 
%of $E_r$.  

%$\bar{q}$ is a member of set $E_r$ and we let $\bar{\theta}_{\bar{q}}=1$ and $\bar{\theta}_q=0$ 
%$\forall q \in E_r \backslash \{\bar{q}\}$.  
%(based on constraint (\ref{conv2}) in (Jaster-DW-1)).

A solution to (Master-DW-1') can easily be mapped to a solution to (SPP-LRS) in a similar way.
For a given  solution $\bar{\theta}$, we can construct $(\bar{T}, \bar{Z})$. Let 
$\bar{q}=(x_{\bar{q}}^T, y_{\bar{q}}^T, t_{\bar{q}}^T, v_{\bar{q}}^T)^T$ be the member 
of $E_r$ such that $\bar{\theta}_{\bar{q}}=1$ and $\bar{\theta}_q=0$ 
$\forall q \in E_r \backslash \{\bar{q}\}$.  
\begin{itemize}
\item $\bar{T}_j=t_{j\bar{q}}$, $\forall j \in J$. %Let $\bar{M}=\{j \in M | \bar{t}_{jq}=1\}$. 
\item Let $\bar{P}_j=\{p_{i_1}, p_{i_2}, .., p_{i_{|\bar{P}_j|}}\} $, $\forall j \in \bar{J}$
be the set of pairings corresponds to vector $x_{\bar{q}}$. 
From proposition \ref{dw-prop-p1} we know that this correspondence is one to one. 
Set $\bar{Z}_p=1$ $\forall p \in \bar{P}_j,j \in \bar{J}$, and 
$\bar{Z}_p=0$ $\forall p \in (\cup_{j \in J \backslash \bar{J}}P_j) \cup (\cup_{j \in \bar{J}} (P_j \backslash \bar{P}_j))$
%$\forall j \in \bar{M}$ and  $\bar{Z}_p=0$ $\forall p \in P_j$ $forall j \in M \backslash \bar{M}$.
\end{itemize} 

Now we show that if $(\bar{T}, \bar{Z})$ is feasible, then $\bar{\theta}=\Psi(\bar{T}, \bar{Z})$
is also feasible. If $(\bar{T}, \bar{Z})$ is feasible, then 
\begin{optprog}
& \sum_{j \in J} \sum_{p \in P_{j}} a_{ip} \bar{Z}_{p} & = &  \sum_{j \in J} \sum_{p \in \bar{P}_{j}} a_{ip}  & = & 1 & \forall i \in N \label{sp_feas} \\
& \sum_{p \in P_j} \sum_{i \in N} a_{ip} D_i \bar{Z}_{p} & = & \sum_{p \in \bar{P}_{j}}\sum_{i \in N} a_{ip} D_i  & \leq & CapF_{j}\bar{T}_j & \forall j \in J \label{fcap_feas} \\
& \sum_{p \in P_j}\bar{Z}_{p} & = & \sum_{p \in \bar{P}_j}\bar{Z}_{p} & \leq &  |H_j| & \forall j \in J \label{numofvec_feas} 
\end{optprog} 

%{\bf(SPP-LRS)} Minimize & \objective{\sum_{j \in M} FC_j T_{j} + \sum_{j \in M} \sum_{p\in P_{j}} C_{p} Z_{p}} \label{spp-obj}\\

The following equations prove that $\bar{\theta}$ is feasible with respect to constraint (\ref{spar}).
\begin{eqnarray}
\sum_{q \in E_r} b_{iq} \bar{\theta}_q  =  b_{i \bar{q}}\bar{\theta}_{\bar{q}}  =  b_{i \bar{q}} =  \sum_{j \in J}g_{ij \bar{q}} =   \sum_{j \in J}\sum_{p \in \bar{P}_j}a_{ip}   =  1  \ \ \forall i \in N  \label{feas-check1} \\
\end{eqnarray} 
Where  $\bar{\theta}_{\bar{q}}=1$ and $\bar{\theta}_q=0$ $\forall q \in E_r \backslash \{\bar{q}\}$.  

The relationship between parameters $g$ and $a$ is from their definition. $g_{ijq}$
is equal to the number of times customer $i$ is served from facility $j$ in solution $q \in E_r$, and $a_{ip}$ 
is equal to 1 if customer $i$ is in pairing $p$ of facility $j$, $\forall i \in N,  p \in P_{j}, j \in J$. Therefore,
to find $g_{ijq}$, we have to sum $a_{ip}$ for all $p \in P_j$ such that $p$ is used to construct solution $q$
$\forall j \in J$, $\forall q \in E_r$. Based
on our notation, $\bar{P}_j$ is the set of pairings used to construct solution $\bar{q}\in E_r$.

We can show that $\bar{\theta}$ is also feasible with respect to constraint (\ref{rel-2}).
\begin{eqnarray}
\sum_{q \in E_r} \sum_{i \in N} g_{ijq} D_i \bar{\theta}_q   & = &   \sum_{i \in N} g_{ij\bar{q}} D_i \bar{\theta}_{\bar{q}}  =   \sum_{i \in N} g_{ij\bar{q}} D_i \nonumber\\
& = & \sum_{i \in N}\sum_{p \in \bar{P}_j}a_{ip}D_i \leq  CapF_{j}\sum_{q \in E_r}t_{jq}\bar{\theta}_q =  CapF_j \bar{T}_j \ \ \forall j \in J \label{feas-check2} 
\end{eqnarray} 

Equations (\ref{feas-check1}) and (\ref{feas-check2}) also
show that the mapped solution  $(\bar{T}, \bar{Z})$ is
feasible when $\bar{\theta}$ is feasible. 
  

Let $z_{SPP}$ be the objective function
value of (SPP-LRS) for solution  $(\bar{T}, \bar{Z})$ and $z_{DW}$ be the objective function
value of (Master-DW-1') for solution  $\bar{\theta}=\Psi(\bar{T}, \bar{Z})$.   
To show that $z_{SPP}= z_{DW}$, we use the definitions of $C$ in (SPP-LRS) and $\tilde{C}$ in (Master-DW-1'). $C_p$ 
is cost of pairing $p$ associated with facility $j$ $\forall p \in P_{j}, j \in J$. This cost includes
travel cost plus the fixed cost of one vehicle. Therefore, for $\bar{\theta}$, if we sum cost of pairings used to 
construct solution $\bar{q} \in E_r$, we can find the total vehicle fixed cost and travel cost  
for variable $\bar{\theta}$
which is $ \sum_{q \in E_r}(VF\sum_{h \in H}v_{hq}+OC \sum_{i \in I} \sum_{k \in I}\sum_{h \in H} TD_{ik}x_{ikhq})\bar{\theta}_q$.
\begin{eqnarray}
z_{DW} = \sum_{q \in E_q}\tilde{C_{q}}\bar{\theta}_q = \tilde{C_{\bar{q}}} & = & \sum_{j \in J}C^F_j t_{j\bar{q}} + \sum_{j \in \bar{J}}\sum_{p \in \bar{P}_j}C_p \nonumber\\
      & = &  \sum_{j \in J}C^F_j \bar{T}_{j} + \sum_{j \in J}\sum_{p \in P_j}C_p Z_p = z_{SPP} \nonumber                       
\end{eqnarray} 
\qed
\end{proof}
\end{comment} 
%%%%%%%%%%%%%%%%%%%%%%%%%%%%%%%%%%%%%%%%%%%%%%%%%%%%%%%%%%%%%%%%%%%%%%%%%%%%%

%%%%%%%%%%%%%%%%%%%%%%%%%%%%%%%%%%%%%%% COMMENT BLOCK %%%%%%%%%%%%%%%%%%%%%%%%%%%%%%%%%%%%%%%%%%%
\begin{comment}
The set of feasible solutions for the system of constraints in $BX'  \leq  D_2$ 
(system (\ref{sys2})) is equal to the set of solutions of the LRS problem, which are 
feasible with respect to the vehicle capacity constraints and the vehicle time 
limit constraints. The system  $AX'  \leq  D_1$ (system (\ref{sys1})) selects
the solution that is feasible with respect to the facility capacity constraints
and set partitioning constraints (\ref{set-part}) among the solutions provided by  
system (\ref{sys2}). 

Instead of system (\ref{sys2}), we can define system $B'X'  \leq  D_2'$, which
generates a set of feasible pairings (with respect to the vehicle capacity constraints 
and the vehicle time limit constraints) for each facility and each vehicle.

 
VERSION 2 : Vehicle index is included to the summation
\begin{optprog}
 $B'X' \leq D_2'= \{$  & \sum_{h \in H_j} \sum_{k \in I} X_{ikh} & \leq & 1 &  \forall  i \in N , j \in J \label{rel-spp}\\
           & \sum_{k\in I} X_{ikh}-\sum_{k \in I} X_{kih} & = & 0 & \forall i \in I, j \in J, h \in H_j \nonumber\\
           & Y_{ikh} - CapV \cdot X_{ikh} & \leq & 0 & \forall i,k \in I, j \in J, h \in H_j\nonumber \\
           & \sum_{k \in I} Y_{ikh} - \sum_{k \in I} Y_{kih} + D_i \sum_{h \in H_j} \sum_{k \in I}X_{ikh} & = & 0 & \forall i \in N, j \in J \label{flowvar} \\
           & \sum_{i \in I}\sum_{k \in I}\left(t_{ik}+s_{k}\right)X_{ikh} - TL\cdot V_h  & \leq & 0 & \forall j \in J, h \in H_j\nonumber\\
           & X_{ikh} & \in & \{0,1\} & \forall i,k \in I, j \in J, h \in H_j\nonumber\\
           & Y_{ikh} & \geq &  0 & \forall i,k \in I, j \in J, h \in H_j \nonumber \\
           & V_h  & \in & \{0,1\} & \forall j \in J, h \in H_j\nonumber\\
\} \nonumber
\end{optprog}

\begin{optprog}
 $B'X' \leq D_2'= \{$  & \sum_{k \in I} X_{ikh} & \leq & 1 &  \forall  i \in N , j \in J, h \in H_j \label{rel-spp}\\
           & \sum_{k\in I} X_{ikh}-\sum_{k \in I} X_{kih} & = & 0 & \forall i \in I, j \in J, h \in H_j \nonumber\\
           & Y_{ikh} - CapV \cdot X_{ikh} & \leq & 0 & \forall i,k \in I, j \in J, h \in H_j\nonumber \\
           & \sum_{k \in I} Y_{ikh} - \sum_{k \in I} Y_{kih} + D_i \sum_{k \in I}X_{ikh} & = & 0 & \forall i \in N, j \in J, h \in H_j \label{flowvar} \\
           & \sum_{i \in I}\sum_{k \in I}\left(t_{ik}+s_{k}\right)X_{ikh} - TL\cdot V_h  & \leq & 0 & \forall j \in J, h \in H_j\nonumber\\
           & X_{ikh} & \in & \{0,1\} & \forall i,k \in I, j \in J, h \in H_j\nonumber\\
           & Y_{ikh} & \geq &  0 & \forall i,k \in I, j \in J, h \in H_j \nonumber \\
           & V_h  & \in & \{0,1\} & \forall j \in J, h \in H_j\nonumber\\
\} \nonumber
\end{optprog}
Where $H_j$ is the set of vehicles located at facility $j$. 
%The constraints  in the system $B'X'  \leq  D_2'$ define the feasible pairings for each facility
%with respect to the vehicle capacity constraints and the vehicle time limit constraints. 
The constraints (\ref{rel-spp}) are relaxed forms of  
constraints (\ref{set-part}). We need the relaxed form because a feasible 
pairing does not need to visit all of the customer nodes but it may not visit a 
customer node more than once. Because of the same reason, in constraint 
(\ref{flowvar}), we multiply $D_i$ with  $\sum_{k \in I}X_{ikh}$.  If the 
pairing does not visit the customer node $i$, then  the flow into the node $i$ 
minus the flow out of the node $i$ must be zero. Otherwise, the summation 
$\sum_{k \in I}X_{ikh}$ is one and the difference between 
the flow into the customer and out of the customer must equal to the demand of 
the customer for pairing assigned to vehicle $h \in H_j$.

Integer feasible points of system (\ref{sys2}) are also integer feasible in 
system $B'X'  \leq  D_2'$ 
(e.g. $\{X' \ \ \mbox{integer}|BX'  \leq  D_2\}\subseteq \{X' \ \ \mbox{integer}|B'X'  \leq  D_2'\}$).
Even the equalities (\ref{flow-conv}) in system (\ref{sys2}) are less strict than 
the equalities (\ref{flowvar}) in system $B'X'  \leq  D_2'$, because of the fact 
that a customer can only be visited by one vehicle, every
integer feasible solution for system  (\ref{sys2}) is also feasible 
for set of inequalities included in system (\ref{sys2}) except
the  equalities (\ref{flow-conv}) are replaced with the equalities (\ref{flowvar}).
Since inequalities (\ref{rel-spp}) are less strict than the less than equal to 
part of equalities (\ref{set-part}) in system (\ref{sys2}), 
$\{X' \ \ \mbox{integer}|BX'  \leq  D_2\}\subseteq \{X' \ \ \mbox{integer}|B'X'  \leq  D_2'\}$
is follows.  

Since each constraint in $B'X'  \leq  D_2'$ is defined for each vehicle 
$h \in H_j$ and facility $j \in J$, we can decompose 
the system $B'X'  \leq  D_2'$ into $|H|$ (where $|H|=\sum_{j \in J}|H_j|$) independent 
subproblems $B'_{h_j}X_{h_j}' \leq D_{2_{h_j}}'$,  one for each vehicle in
each facility. Each subproblem 
generates feasible pairings for a vehicle located at a facility. We define  
$P_{h_j}$ as the set of feasible pairings for vehicle $h$ at facility $j$.
However, since the vehicles in facility $j$ are identical, we know that
the sets $P_{h_j}$ are identical $\forall h \in H_j$. Therefore, instead of sets 
$P_{h_j}$ $\forall h \in H_j$, we can use the set $P_j$, the set of feasible pairings 
for facility $j$. If we define  $x_{ikp}^{j}$ as the number of times arc $(i,k)$ used in pairing $p$
(resource constrained path that originates from facility $j$, visits some set of 
customers and facility $j$ and returns to facility $j$), then we can represent the vector $X$ in  
master problem as a nonnegative linear combination of a finite number of pairings:

\begin{optprog}
& X_{ikh_j} & = & \sum_{p \in P_{j}}x_{ikp}^{j} \cdot \theta_p^{h_j} &  \forall  i,k \in I , j \in J, h_j \in H_j \label{XintermsTeta}\\
& \sum_{p \in P_{j}}\theta_p^{h_j} & \leq & 1 & \forall j \in J, h_j \in H_j \\
&  \theta_p^{h_j}  & \in & \{0,1\} & \forall j \in J, h_j \in H_j, p \in P_{j}
\end{optprog}  
where 
\begin{eqnarray} 
\theta_p^{h_j} = \mbox{number of times pairing $p$ of vehicle $h_j$ is selected. }
\end{eqnarray}
If we sum flow conservation constraints (\ref{flowvar}) over all $i \in N$, simplify appropriate
parts, use the information that the vehicle turns back to the depot with zero load because of the characteristics 
of equations (\ref{flowvar}) and substitute equality (\ref{XintermsTeta}),  we get
\begin{eqnarray} 
\sum_{i \in N} Y_{jih_j} =  \sum_{i \in N}D_i \sum_{k \in I} \sum_{p \in P_{j}}x_{ikp}^{j} \cdot \theta_p^{h_j} \ \ \forall j \in J, h_j \in H_j \label{YintermsofTeta}
\end{eqnarray} 
To transform variable $V_h$ in master problem, we use the following equality:
\begin{eqnarray}
V_{h_j}= \sum_{p \in P_j} \theta_p^{h_j} \ \ \forall h_j \in H_j, j \in J  \label{driintermsTeta}
\end{eqnarray} 
Subproblems for $h_j \in H_j$ and facility $j$ are independent and  identical. Therefore, 
%instead of $P_{h_j}$, we can define $P_j$ as the set of pairings for facility $j$ and 
we can introduce variable $\theta_p^{j}$ satisfying following system of constraints.
\begin{optprog}
& \theta_p^{j} & = & \sum_{h_j \in H_j} \theta_p^{h_j} & \forall p \in P_j \label{aggteta}\\ 
& \sum_{p \in P_j} \theta_p^{j} & \leq & |H_j|  & \forall j \in J \\
& \theta_p^{j}  & \in & \{0,1\} & \forall j \in J, p \in P_j  \label{teta_jbinary}
\end{optprog}  


Constraints (\ref{rel-1}) are used to represent the constraints 
(\ref{Rel-X-T-1}). We obtain constraints (\ref{rel-1}) 
using the following relations.
\begin{optprog}
& X_{ijh_j} & = & \sum_{h_j \in H_j} X_{ijh_j} &  \forall  i \in N , j \in J, h_j \in H_j \label{first-rel}\\
&\sum_{h_j \in H_j} X_{ijh_j} & = & \sum_{p \in P_j} x_{ijp}^{j} \cdot \theta_p^{j} &  \forall  i \in N , j \in J \\
& \sum_{p \in P_j} x_{ijp}^{j} \cdot \theta_p^{j} & \leq & \sum_{p \in P_j}a_{ip}\theta^j_p  \leq 1  &  \forall  i \in N , j \in J  \\
& \sum_{p \in P_j}a_{ip} \theta^j_p & \leq & T_{j}  & \forall j \in J,  \forall i\in N \label{relation_teta}  
\end{optprog} 
The last relations (\ref{relation_teta}) are the result of constraints (\ref{rel-2}). Relations
from (\ref{first-rel}) to (\ref{relation_teta})  give  inequalities (\ref{Rel-X-T-1}), which are 
$X_{ijh_j}\leq T_j$ $\forall  i \in N , j \in J, h_j \in H_j$. 

\end{comment}

%%%%%%%%%%%%%%%%%%%%%%%%%%%%%%%%%%%%%%%%%%%%%%%%%%%%%%%%%%%%%%%%%%%%%%%%%%%%%%%%%%%%%%%%%%%%%%%%%%
%%%%%%%%%%%%%%%%%%%%%%%%%%%%%%%%%%%%%%% COMMENT BLOCK %%%%%%%%%%%%%%%%%%%%%%%%%%%%%%%%%%%%%%%%%%%
%%%%%%%%%%%%%%%%%%%%%%%%%%%%%%%%%%%%%%%%%%%%%%%%%%%%%%%%%%%%%%%%%%%%%%%%%%%%%%%%%%%%%%%%%%%%
% rewrite this part    
%Constraints (\ref{rel-1}) are implied by Constraints (\ref{rel-2}). However, constraints (\ref{rel-1})
%are valid inequalities, thus we can add these constraints that are representing the relationship 
%between variable $T_j$ and $\theta^j_{p}$ to our set partitioning-based formulation SPP-LRS. 
%In SPP-LRS, those constraints will be 
%$\sum_{p \in P_j}a_{ip}Z_{p} \leq  T_{j}\ \ \forall j \in M, \forall i \in N$. 
%$Z_{p}  \leq  T_{j}\ \ \forall j \in M, \forall p \in P_{j}$. 



% and $a_{pj}$ as the column 
%vector such that  
%\begin{eqnarray*}
%a_{ip} & = & \left \{  \begin{array}{ll}
%                 1 & \mbox{if node $i$ is in pairing $p$ of facility $j$, $\forall i \in N, j \in M, p \in P_{j}$ } \\
%                 0 & \mbox{otherwise}
%                         \end{array}
%              \right. \\
%C_{p} &=&\mbox{cost of the pairing pairing $p$  of facility $j$ (column $a_{pj}$), $\forall j \in M, p \in P_{j}$} 
%\end{eqnarray*} 
%Since we are solving one subproblem for each facility, we define our node set 
%as $I^j=N \cup \{j\}$. We will create one feasible pairing, so we can drop the 
%index h from the original variable definition and define variables 
%$(X_{ik},Y_{ik})^{N+1,N+1}$:
%\begin{eqnarray*}
%X_{ik} & = & \left\{ \begin{array}{ll}
%                 1  & \mbox{if arc (i,k) is selected $\forall i \in I^j, k \in I^j$}\\
%                 0  & \mbox{otherwise}                                                  
%                \end{array} 
%              \right. \\
%Y_{ik} & = & \mbox{flow from location $i$ to location $k$, $\forall i \in I^j, k \in I^j$}
%\end{eqnarray*}
%Each $(X_{ik},Y_{ik})^{N+1,N+1}$ vector represents one feasible pairing for facility $j$.
%We can define the set $P_j$ by using the system $B'_jX_j' \leq D_{2j}'$:

%\begin{optprog}
%$P_j= \{(X_{ik},Y_{ik})^{N+1,N+1}:$  & \sum_{k \in I^j} X_{ik} & \leq & 1 &  \forall  i \in N \nonumber \\
%           & \sum_{k\in I^j} X_{ik}-\sum_{k \in I^j} X_{ki} & = & 0 & \forall i \in I^j \nonumber\\
%           & Y_{ik} - CapV \cdot X_{ik} & \leq & 0 & \forall i,k \in I^j \nonumber \\
%           & \sum_{k \in I^j} Y_{ik} - \sum_{k \in I^j} Y_{ki} + D_i \sum_{k \in I^j}X_{ik} & = & 0 & \forall i \in N \label{f:con} \\
%           & \sum_{i,k \in I^j}\left(t_{ik}+s_{k}\right)X_{ik} - TL  & \leq & 0 \nonumber\\
%           & X_{ik} & \in & \{0,1\} & \forall i,k \in I^j \nonumber\\
%           & Y_{ik} & \geq &  0 & \forall i,k \in I^j \nonumber
%\}
%\end{optprog}

%The first constraint in the definition of $P_j$ is the relaxed form of our set 
%partititoning constraint (\ref{set-part}). We have to use the relaxed form because in a 
%feasible pairing it is not necessary to visit all customer nodes, but still a 
%customer node can be visited at most once. 
 
%Define 
%\begin{eqnarray*}
% \lambda^j_p & = & \mbox{the number of times pairing $p$ associated with facility $j$ (column $a_{pj}$) is}\\ 
%             & & \mbox{selected in the solution of (EB-LRS'),} \ \ \forall j \in M, p \in P_j
%\end{eqnarray*}

%To construct the Dantzig Wolfe Decomposition of the problem (EB-LRS'), we write $(X,Y)^{pj}$, 
%the vector of variables defining pairing $p$ for facility $j$, in terms of column $a_{pj}$ and
%variable $\lambda^j_p$. Let $X^{B'X'\leq D_2'}$ define all feasible points satisfying 
%the system $B'X'\leq D_2'$. Then we have the following relations: 

%Old version
%\begin{optprog}
%$X^{BX\leq D_2}$ &= \{(X,Y)^j_{j \in M} : (X,Y)^j &\in & P_j \ \ \forall j \in M \} \nonumber\\
%                 &= \{(X,Y)^j_{j \in M} : X^j_{ik} &\leq & a_{ip}\lambda^j_p \leq \lambda^j_p & \forall  i,k \in I^j, j\in M, p \in P_j \nonumber \\
%                 & \sum_{k \in I^j} X^j_{ik} &=& a_{ip}\lambda^j_p  & \forall  i \in I^j, j\in M, p \in P_j \nonumber \\
%                  & -\sum_{k \in I^j} Y^j_{ki} &=& -\sum_{i \in N}D_i a_{ip}\lambda^j_p  & \forall j\in M, p \in P_j  \label{f:con:tran}\\ 
%                 & \sum_{p \in P_j}\lambda^j_p &\leq & |H^j|  & \forall j \in M \label{conv}\\
%                 & \lambda^j_p &\in & \{0,1\} & \forall j \in M, p \in P_j \}\nonumber  
%\end{optprog}

%\begin{optprog}
%$X^{B'X\leq D_2'}$ &= \{(X,Y)^j_{j \in M} : (X,Y)^j &\in & P_j \ \ \forall j \in M \} \nonumber\\
%                 &= \{(X,Y)^j_{j \in M} : X^j_{jk} &\leq & \sum_{p \in P_j}a_{kpj}\lambda^j_p  & \forall  k \in N, j\in M \label{forXTrel1}\\
%                 & X^j_{kj} &\leq & \sum_{p \in P_j}a_{kpj}\lambda^j_p  & \forall  k \in N, j\in M \label{forXTrel2}\\
%                 & \sum_{k \in I^j} X^j_{ik} &=& \sum_{p \in P_j}a_{ip}\lambda^j_p  & \forall  i \in N, j\in M  \label{forsetpart} \\
%                  & -\sum_{i \in N} Y^j_{ji} &=& -\sum_{i \in N}\sum_{p \in P_j}D_i a_{ip}\lambda^j_p  & \forall j\in M  \label{f:con:tran}\\ 
%                 & \sum_{p \in P_j}\lambda^j_p &\leq & |H^j|  & \forall j \in M \label{conv}\\
%                 & \lambda^j_p &\in & \{0,1\} & \forall j \in M, p \in P_j \}\nonumber  
%\end{optprog}

%[{\bf Explanation:}
%-We could not give the equivalence of variables $X^j_{ik}$ and $Y^j_{ik}$ (e.g., $X^j_{ik}=?$ and $Y^j_{ik}=?$). But to represent 
%system (\ref{sys1}) in terms of variables $\lambda$, we need relations (\ref{forXTrel1}), (\ref{forXTrel2}), (\ref{forsetpart})
%and (\ref{f:con:tran}). 

%-Relations (\ref{forXTrel1}) and (\ref{forXTrel2}): We need them to represent Const.(\ref{Rel-X-T-1}), Const.(\ref{Rel-X-T-2}). 
%They are true, since if $X^j_{jk}=0$ then $\sum_{p \in P_j}a_{kpj}\lambda^j_p$ may be 0 or 1
%and if $X^j_{jk}=1$ then $\sum_{p \in P_j}a_{kpj}\lambda^j_p$ will be 1. Also, to write constraints (\ref{Rel-X-T-1}) and (\ref{Rel-X-T-2}), 
%we need $\sum_{p \in P_j}a_{kpj}\lambda^j_p \leq T_j$.          
%Based on this new relation, the constraint (\ref{rel-1}) should be updated to
%$\sum_{p \in P_j}a_{ipj}\lambda^j_p  \leq  T_{j} \forall j \in M, \forall i\in N $
%-Relations (\ref{forsetpart}): We need them to represent Const.(\ref{set-part}). 
%-Relations (\ref{f:con:tran}): We need them to represent Const.(\ref{FacCap}).

%\begin{figure}[h]
%  \centering
%  \includegraphics[scale=1]{figs/danzigfig.eps}
%\end{figure} 

%For example based on the figure above, we can write flow conservation equation (\ref{f:con}) for all $i \in N$:
%\begin{eqnarray*}
%(Y_{13}+Y_{14}+Y_{12}+Y_{15}+Y_{10})-(Y_{01}+Y_{31}+Y_{21}+Y_{51}+Y_{41})+D_1\sum{k \in I^j}X_{1k}=0 \\
%(Y_{31}+Y_{34}+Y_{32}+Y_{35}+Y_{30})-(Y_{03}+Y_{13}+Y_{23}+Y_{53}+Y_{43})+D_3\sum{k \in I^j}X_{3k}=0 \\
%(Y_{43}+Y_{41}+Y_{42}+Y_{45}+Y_{40})-(Y_{04}+Y_{34}+Y_{24}+Y_{54}+Y_{14})+D_4\sum{k \in I^j}X_{4k}=0 \\
%(Y_{23}+Y_{24}+Y_{21}+Y_{25}+Y_{20})-(Y_{02}+Y_{32}+Y_{12}+Y_{52}+Y_{42})+D_2\sum{k \in I^j}X_{2k}=0 \\
%(Y_{53}+Y_{54}+Y_{52}+Y_{51}+Y_{50})-(Y_{05}+Y_{35}+Y_{25}+Y_{15}+Y_{45})+D_5\sum{k \in I^j}X_{5k}=0 
%\end{eqnarray*}
%When we sum equation (\ref{f:con}) for all $i \in N$, we get $\sum_{i \in N}Y_{i0}-\sum_{i \in N}Y_{0i}+\sum_{i \in N}D_i \sum_{k \in I^j}X_{ik}=0$.
%We know that $\sum_{i \in N}Y_{i0}=0$. 
%]

 
%Equality (\ref{f:con:tran}) is obtained by adding the flow conservation equation (\ref{f:con}) for all $i \in N$ and
%using the information that the vehicle turns back to depot with zero load because of the characteristics 
%of equations (\ref{f:con}) in system (\ref{sys2}).    

%By using the transformation of $(X,Y)^{pj}$ to columns $a_{pj}$ and variable $\lambda^j_p$, we can write the
%master problem, which is the system (\ref{sys1}), $AX'\leq D_1$, with the objective function. To write the objective 
%function, we use the cost definition of the pairing, that is $C_{p}=VF+OC \sum_{i,k \in I^j}(t_{ik}+s_k)X_{ik} $. 

%\begin{optprog}
%{\bf (Master-DW)} \nonumber\\
%Minimize & \objective{\sum_{j \in M} FC_j T_{j} + \sum_{j \in M} \sum_{p\in P_{j}} C_{p} \lambda^j_p} \nonumber\\
%subject to & \sum_{j \in M} \sum_{p \in P_{j}} a_{ipj} \lambda^j_p & = & 1 & \forall i \in N \label{spar} \\
%           & \sum_{p \in P_j} \sum_{i \in N} a_{ipj} D_i \lambda^j_p  & \leq & CapF_{j} T_{j} & \forall j \in M \label{rel-2}  \\
%           & \sum_{p \in P_j}a_{ipj}\lambda^j_p & \leq & T_{j}  & \forall j \in M,  \forall i\in N \label{rel-1} \\
%           & \sum_{p \in P_j} \lambda^j_p  & \leq & H_j & \forall j \in M \label{conv2} \\
%           & \lambda^j_p  & \in & \{0,1\} & \forall  j \in M, \forall p \in P_{j} \label {conv3}\\
%           & T_{j}  & \in & \{0,1\} & \forall  j \in M \nonumber
%\end{optprog} 

%To write Master-DW in terms of new variables, along with the definition of $X^{B'X'\leq D_2'}$, we
%have used some logical equalities. For example, in constraint (\ref{spar}) we have used the 
%transformation $\sum_{h \in H} \to \sum_{j \in M} \sum_{p \in P_{j}}$ along with the equality 
%$\sum_{k \in I^j} X^j_{ik} = a_{ipj}\lambda^j_p $. 


%%%%%%%%%%%%%%%%%%%%%%%%%%%%%%%%%%%%%%%%%%%%%%%%%%%%%%%%%%%%%%%%%%%%%%%%%%%%%%%%%%%%%%%%%
%%%%%%%%%%%%%%%%%%%%%%%%%%%%%%%%%%%%%%%%%%%%%%%%%%%%%%%%%%%%%%%%%%%%%%%%%%%%%%%%%%%%%%%

\section{Comparing G-LRSP and SPP-LRSP}
\label{sec:comp-form}
% Number of variables and constraints
G-LRSP is a three-index commodity flow formulation.
For the related problems of LRP and VRP, the literature contains 
two main types of graph-based models, vehicle flow and 
commodity flow. In both formulations, integer variables 
represent the number of times a vehicle traverses a 
given arc or edge in an underlying graph. In commodity
flow formulations, an additional variable representing the 
number of units of a given commodity transported along the 
arc is present. In general, vehicle flow models contain an 
exponential number of subtour elimination constraints, 
whereas commodity flow models, with the help of additional 
flow variables, use a polynomial number of constraints to 
eliminate subtours. For the LRP and the VRP, two-index 
vehicle flow formulations are preferred in solution algorithms 
based on branch-and-cut since these formulations include 
a smaller number of variables and their LP relaxations 
yield better bounds. However, it is not possible to formulate 
the LRSP using variables with only two indices because of 
the time constraint for the vehicles. In this context, 
we require a three-index commodity flow model.
%G-LRSP is a three-index commodity flow formulation. 
Because of the facility and vehicle capacity restrictions, 
we do not know a priori the number of vehicles required at 
each facility, though the number is clearly at most $|I|$. 
Thus, this formulation can require up to O($|I|^4$) binary 
variables and O($|I|^4$) constraints, which can be as many as 
6 million for an instance with 50 customers. 

On the other hand, SPP-LRSP includes O($|I|$) number of constraints
while the number of variables depends on the size of the pairing 
sets (i.e., $|\bigcup_{j \in J}P_j|$), which can be exponentially 
large. However, most of these columns are equal to zero in an 
optimal solution. Therefore, intelligent enumeration techniques 
which may use some logical conditions, information from good 
solutions, and information from experiences can be developed to 
avoid generating all possible columns. In addition, the 
disadvantage of having extremely large number of variables can 
be overcome by using column generation algorithm that explores
all columns implicitly and generates only those that may improve 
the objective function. 
%By solving a 
%sequence of linear programs and dynamically generating columns 
%eligible to enter the basis, such an algorithm implicitly 
%considers all columns but generates only those that may improve 
%the objective function. 
  

%Symmetry
Observe that the three-index formulation G-LRSP exhibits a high 
degree of symmetry under the assumption that the vehicle 
fleet assigned to each facility is homogeneous. This is due to the 
fact that the assignment of pairings to a specific vehicle is
essentially arbitrary, i.e., the cost of a given solution to 
G-LRSP is invariant under permutation of the indices assigned to 
specific vehicles. In addition, from the definition of the variables
in formulation G-LRSP, symmetric distance matrix results in solutions
with same cost but including routes constructed in different 
directions. A solution algorithm such as a branch-and-bound
or a branch-and-cut based on G-LRSP may perform poorly
because of the inherent symmetry in the formulation. 

The SPP-LRSP formulation eliminates the symmetry 
that arises in G-LRSP because of the identical vehicles. In SPP-LRSP,
we enumerate one set of feasible pairings that can be assigned to any 
vehicle located at a facility instead of duplicating this set for 
each vehicle located at a facility. Furthermore, while generating
the pairings using a subproblem, depending on the solution procedure, 
we can use undirected links between nodes in case of symmetric distance
matrix. Therefore, we can remove direction in the routes.  

The lower bound obtained by solving the  LP relaxation of
the SPP-LRSP is stronger than the bound obtained by solving the  
LP relaxation of the G-LRSP. This can be demonstrated both theoretically
and empirically. Let $z_{LP}^{G}$ be the LP relaxation of formulation
G-LRSP and $z_{LP}^{SPP}$ be the LP relaxation of formulation
SPP-LRSP. In Proposition \ref{dw-prop-p2} we proved that 
MDW' obtained by applying Dantzig Wolfe Decomposition to G-LRSP
is equivalent to SPP-LRSP. Therefore using DWD methodology we 
can theoretically compare the LP relaxation bounds of G-LRSP 
and SPP-LRSP. Let $z_{LP}^{MDW'}$ be the LP relaxation of
formulation MDW'. Then, we have the following result. 
\begin{prop}
\label{ch1-comp-bnds}
$z_{LP}^{G}\leq z_{LP}^{MDW'} = z_{LP}^{SPP}$.
\end{prop}
\begin{proof}
$z_{LP}^{G}$ is obtained by optimizing the objective function
(\ref{obj-eb}) on the polyhedron defined by constraints 
(\ref{set-part}) - (\ref{fix-0}) (G-LRSP without the integrality 
constraints). 

On the other hand, $z_{LP}^{MDW'}$
is obtained by optimizing the objective function (\ref{obj-eb})
on the polyhedron defined by the intersection of constraints 
(\ref{set-part}) and (\ref{FacCap}) with the convex hull of
of the solutions that are constructed with composing the points 
generated by solving subproblem (\ref{subsys}) for all facilities,
which is the convex hull defined by (\ref{enterleave}), 
 (\ref{VehCap}), (\ref{flow-conv}), (\ref{timeLim}), (\ref{fix-0}), 
(\ref{x-var}), (\ref{y-var}), (\ref{t-var}), and (\ref{dri-var}).

The proof follows from the fact that the subproblems do not have 
the integrality property, i.e., the solution is unchanged when the 
integrality restriction is removed. The polyhedron defined by the 
constraints (\ref{enterleave}), (\ref{VehCap}), (\ref{flow-conv}), 
(\ref{timeLim}), (\ref{fix-0}) is greater or equal to the 
convex hull of the same set of constraints.  
\qed
\end{proof} 

The relation in Proposition  \ref{ch1-comp-bnds} can also
be demonstrated emprically. In Section \zaupdate{REFER COMPUTATIONAL
RES. FOR LP RELAX.}, 
we present computational results that compare the LP relaxation bounds 
obtained by solving SPP-LRSP and G-LRSP.  
 

Finally, while we are searching for cutting planes valid for the model, 
we can see some differences between formulation G-LRSP and formulation 
SPP-LRSP. In formulation G-LRSP, the variables indexed with vehicle 
$h$ which enables us to search cutting planes associated with either a single 
vehicle or a set of vehicles. We can multiply a cut generated for 
a single vehicle $h$ for all vehicles in set $H$. This can prevent jumps between 
the symmetrical solutions because of the homogeneous vehicle fleet. 
However, in formulation SPP-LRSP, a pairing variable for a facility
is associated with any vehicle located at the facility, thus we can
not differentiate the vehicles. We can only use cutting planes 
associated with all vehicles located at a facility. These cutting
planes have the same coefficient for all pairings (representing
a single vehicle) located at a facility. 
Solution process for SPP-LRSP becomes more complicated if cutting
planes only including a subset of pairings, and possibly with 
different coefficients for the variables representing pairings
are added. Therefore, the structure of the cutting planes searched 
is more restrictive for SPP-LRSP. However, one advantage
of SPP-LRSP is that some cutting planes valid for formulation G-LRSP
are already implied by formulation  SPP-LRSP. This is evident from
the Proposition  \ref{ch1-comp-bnds} and the definition 
of the columns of SPP-LRSP. This will further be investigated in 
Section \zaupdate{REFER CUTS SECTION}.



%Look at Barnhart's paper, why need models with large number of variables..
%BP 
   

%1st paragraph: LP relaxation bound differences
%2nd paragraph: EB is very symmetric. Symmetric wrt routes, e.g. 
%forward and backward is same. Symmetric wrt to vehicles. Any vehicle 
%can be assigned to any route. With this symmetry branch and cut is almost 
%hopeless in EB-LRS.
%In SP-LRS, this symmetry is eliminated. Symmetry wrt to edges is eliminated. 
%Also a column represent a group of vehicles so eliminates the symmetry because 
%of identical vehicles. Which is a good thing. 

%A cut in edge based can be multiplied for the other vehicles. 
%But a cut for two vehicles cannot be written in sp model. 










\section{Complexity of the LRSP}
The LRSP is NP-hard. To show this, we consider an instance of the LRSP 
in which there is at most one vehicle at each facility, there are 
no facility and vehicle fixed costs, and the vehicle time limit is 
unrestricted. Such a special case is equivalent to an instance of 
the MDVRP, which was shown to be NP-hard by  \citet{Lenstra},
and \citet{Bodin2}.
%To show that the LRSP is NP-hard, we start from the fact that the
%multi-depot vehicle routing problem (MDVRP) is NP-hard 
%(\citet{Bodin2,Lenstra}). 
The MDVRP aims to find a minimum-cost set of vehicle routes, each 
one starting from a depot and terminating at the same depot. Each 
customer must be served exactly once and there are specified lower and 
upper bounds for the number of capacitated vehicles at each depot. 
%To show that the LRSP is NP-hard, we show that the 
%MDVRP is a special case of LRSP.  
%We will consider an instance of the LRS 
%problem with $N$ customer nodes and $J$ facilities and at most $I$ vehicles at 
%each depot. 
In the created LRSP instance, location decision and facility costs are 
eliminated by setting the fixed cost of facilities to zero, and scheduling
part is eliminated by setting the time limit to infinity. Solution of this
instance results a set of routes with minimum total operating 
cost.
%Since the MDVRP considers only the operating costs of vehicles, the 
%fixed cost of a vehicle (in the LRS problem) is assumed 
%to be zero for the instance. 
%The total working hour limit of a vehicle is 
%eliminated, or equivalently the total working hour limit of the vehicle is set to infinity.
%The created instance of the  LRS problem is equivalent to an instance of the MDVRP 
%with $I$ customer nodes and $J$ facilities, each of which has $0$ to $I$ vehicles.
This solution can easily be translated to a solution for the MDVRP. 
In the solution of the MDVRP, the allocation of the customers to the facilities, 
the vehicle routes and the total cost are the same as in the solution of the LRSP. 
%The only difference is the number of vehicles used in the solution. 
%In the MDVRP, each route is assigned to one vehicle, therefore
%the number of vehicles used is equal to the number of routes. 
%However, in the created instance of the LRS problem, there is one vehicle at each facility and 
%all the routes originated from a facility are assigned to the vehicle located at 
%the facility. 
Since we have shown that the MDVRP is a special case of the LRSP and the 
MDVRP is NP-hard, we can conclude that the LRSP is also NP-hard. 
 



%If both problems have the same distance matrix and same operating cost, then solving the LRS 
%problem yields a solution for the MDVRP.
 
%Each facility has a capacity of value $FCap$ and has 1 vehicle with capacity $VCap$ and no total working hour 
%limit, which means that we set the 
%total working hour limit of each vehicle to infinity.  This instance for 
%LRS problem can be converted to an instance of MDVRP with M facilities 
%and N customer nodes in which each facility has 0 to 
%N vehicles with capacity $VCap$. If both problems have the same 
%distance matrix and same operating cost, then solving the LRS problem 
%is equavalent to solving the MDVRP. 

%We define a special case of LRS problem,
%solution of which is equivalent to solution of MDVRP.
%%%%%%%%%%%%%%%%%%%%%%%%%%%%%%%%%%%%%%%%%%%%%%%%%%%%%%%%%%%%%%%%%%%
%[Dr. Berger deleted this part] ??
%The allocation of customers
%to facilities, vehicle routes and the total costs are same. 
%The only difference is in MDVRP at most one route is assigned to
%each vehicle although in LRS problem it is possible to assign
%multiple routes to a vehicle. Thus, if multiple vehicles are used to
%serve customers in any depot, while converting the 
%solution of MDVRP to the solution of LRS, 1 vehicle is assigned to these 
%multiple routes in that depot. This assignment is possible since there is no working hour 
%limit for the vehicles. Since a special case of LRS problem is equivalent to
%a MDVRP, LRS problem is NP-Hard.           






%%%%%%%%%%%%%%%%%%%%%%%%%%%%%%%%%%%%%%%%%%%%%%%%%%%%%%%%%%%%%%%%%%%%%%%%%%%%%%%%%%%%%%%%%%%%%%%%%
\begin{comment}
{\bf LP Relaxation of Master Problem:} In this part, we use (Master-DW-2)
as our master problem obtained by applying Dantzig Wolfe Decomposition 
method to the EB-LRS. The subproblem, the resource constrained
elementary shortest path problem does not have integrality property. Therefore, 
the bound obtained from Dantzig Wolfe Decomposition of LRSP is expected to be equal or
stronger than the LP relaxation of EB-LRS. Let $z_{LP}^{EB-LRS}$ be the LP
relaxation of model EB-LRS and $z_{LP}^{Master-DW}$ be the bound obtained from the Dantzig 
Wolfe Decomposition of EB-LRS. Then,
\begin{prop}
$z_{LP}^{EB-LRS}\leq z_{LP}^{Master-DW}$.
\end{prop}
\begin{proof}
Let $Q^A$ be the polyhedron defined by $AX' \leq D_1$, i.e 
$Q^A=\{X'\in \mathbb{R}^{|D|}_+|AX' \leq D_1\}$. Let $Q^B$ be the polyhedron 
defined by $BX' \leq D_2$, i.e.,, $Q^B=\{X'\in \mathbb{R}^{|D|}_+|BX' \leq D_2\}$.
The LP relaxation of the edge based formulation (EB-LRS) is
\begin{eqnarray}
\label{lpbound-eb}
z_{LP}^{EB}=min \{CX' | X' \in (Q^A \cap Q^B),X'\in \mathbb{R}^{|D|}_+\} \nonumber
\end{eqnarray}

We decompose the subproblem into $|M|$ subproblems, one for each facility. 
Let $B_jX' \leq D_{2j}$ be the system of inequalities for subproblem $j$.
Let $S^{B}$ be the set of integer points satisfying $BX' \leq D_2$.
\begin{eqnarray}
S^{B} & = & \{X'\in \mathbb{Z}^{|D|}_+|BX' \leq D_2\} \nonumber\\
      & = & \bigcup_{j \in J} (\{X'\in \mathbb{Z}^{|D^j|}_+|B_jX' \leq D_{2j}\} \nonumber 
\end{eqnarray}
Let $P^{B}$ be the convex hull of the points in set $S^{B}$. 
We know that $Q^{B}\supseteq P^{B}$. Therefore, 
$(Q^A \cap Q^B) \supseteq (Q^A \cap P^{B})$, which proves the proposition.
\qed
\end{proof}

Dantzig Wolfe Decomposition bound $z_{LP}^{Jaster-DW}$ can be calculated using 
a column generation algorithm. The details of the column generation algorithm is
explained in part II, section (\ref{ColGenAlg}). The process starts with initializing the
(Master-DW) with some set of columns. The LP relaxation of (Master-DW) with the
current set of columns is optimized, then using the dual variable information
the subproblem is solved to find minimum reduced cost columns. 

Let $\pi_{i}$ be the dual variable 
associated with the partitioning constraint (\ref{spar2}) for customer $i$, $\mu_{j}$ be 
the dual variable associated with the capacity constraint (\ref{rel-2-2}) for 
facility $j$, $v_j$ be the dual variable associated with the convexification
constraint  (\ref{up-bound-lambda}) for facility $j$. Based on this information, we
can write the reduced cost of a column which will lead to the objective function
coefficients of the subproblem (Dec-Sub-Prob-2). Let $\hat{C}_{p}$ be the reduced
cost of column $p$ such that $p \in P_j$.
\begin{eqnarray}
\label{dw-red-cost}
\hat{C}_{p}=C_{p}-\sum_{i\in N}a_{ip}\pi_{i}+\sum_{i \in N}a_{ip}D_{i}\mu_{j}+v_{j}
\end{eqnarray}
$C_{p}$ is the cost of a column which is the sum of the fixed cost of the vehicle and 
the variable costs of the routes that are included in the pairing.
Let $x_{ikp}$ be the binary parameter which is equal to 1 if pairing $p$ 
visits link $(i,k)$ and 0 otherwise.
Using parameter $x$ we can rewrite (\ref{dw-red-cost}).
%To drive the objective function coefficients in subproblem (Dec-Sub-Prob-2) for
%facility $j$, we substitute equalities (\ref{transform-xa}) and (\ref{transform-cost}):
\begin{eqnarray}
\hat{C}_{p} & = & VF+OC \sum_{i \in I} \sum_{k \in I} t_{ik}x_{ikp}-\sum_{i\in N} \pi_{i} \sum_{k \in I}x_{kip}+ \mu_{j}\sum_{i \in N}D_{i}\sum_{k \in I}x_{kip}+v_{j} \\
            & = & (VF + v_j)+ \sum_{k \in I}\sum_{i \in N}(OC t_{ki} - \pi_{i} + \mu_{j}D_{i})x_{kip} + OC\sum_{i \in N}t_{ij}x_{ijp} \label{openform-rc}
\end{eqnarray} 
Note that in subproblem for facility $j$, set of all nodes $I$ is equal to customer 
nodes plus the facility $j$, i.e., $I=N\cup\{j\}$. Based on equation (\ref{openform-rc}), 
for objective function (\ref{de-subprob-obj}), we can define:
\begin{eqnarray*}
\hat{c}_{ki} & = & \left \{  \begin{array}{ll}
                 (OC t_{ki} - \pi_{i} + \mu_{j}D_{i})  & \mbox{if $k \in I$, and $i \in N$} \\
                 (OC t_{ki}) & \mbox{if $k \in N$ and $i=j$}
                         \end{array}
              \right. \\
K & = & VF + v_j
\end{eqnarray*}
\end{comment}
%%%%%%%%%%%%%%%%%%%%%%%%%%%%%%%%%%%%%%%%%%%%%%%%%%%%%%%%%%%%%%%%%%%%%%%%%%%%%%%%%%%%%%%%%%%%%%%%%







